\documentclass{article}
% NOTE: the NeurIPS 2027 style file is not yet distributed. We use the 2026
% style as a placeholder and will switch on release.
\usepackage[preprint,nonatbib]{neurips_2026}
\usepackage[utf8]{inputenc}
\usepackage[T1]{fontenc}
\usepackage{hyperref}
\usepackage{url}
\usepackage{booktabs}
\usepackage{graphicx}
\usepackage{amsfonts}
\usepackage{array}
\usepackage{tabularx}
\usepackage{float}
\usepackage{microtype}
\usepackage{xcolor}
\usepackage{enumitem}
\usepackage{mdframed}
\usepackage{amsmath}

\newenvironment{promptbox}{\begin{mdframed}[linewidth=0.4pt,linecolor=black,innerleftmargin=5pt,innerrightmargin=5pt,innertopmargin=5pt,innerbottommargin=5pt,skipabove=2pt,skipbelow=5pt]\setlength{\parskip}{0pt}\setlength{\parindent}{0pt}\ttfamily\scriptsize\raggedright\frenchspacing\obeylines\leavevmode\strut\ignorespaces}{\end{mdframed}}

% ---------------------------------------------------------------------------
% Result constants.  Every number in the paper is defined once here so the
% text, tables, and abstract cannot drift out of sync.
% ---------------------------------------------------------------------------
\newcommand{\FASTERBase}{0.93}          % FASTER median Mops/s, same box/harness
\newcommand{\HydraFinal}{5.51}          % hardened (task-3) engine, original workload
\newcommand{\HydraFinalRuns}{5.51 / 5.51 / 5.53}
\newcommand{\HydraProd}{5.52}           % production (task-4) engine, original workload
\newcommand{\HydraProdRuns}{5.50 / 5.51 / 5.53 / 5.51 / 5.53 / 5.53}
\newcommand{\SpeedupFinal}{5.9$\times$}
\newcommand{\HydraOrigVFourC}{5.48}     % v4c engine, original workload
\newcommand{\HydraSparse}{5.67}
\newcommand{\HydraClustered}{5.56}
\newcommand{\HydraDrift}{3.97}
\newcommand{\HydraHoldout}{3.87}
\newcommand{\HoldoutRuns}{3.87 / 3.87 / 3.21}
\newcommand{\HydraFirst}{1.85}          % first (dense-index) engine
\newcommand{\HydraAuditVerified}{5.51}  % July-20 independent audit re-run, production engine
\newcommand{\HydraAuditFixed}{5.50}     % audit-fix (task-5) engine, same-day medians (A/B control also 5.50)
\newcommand{\HydraFinalWL}{5.51}        % final (task-6, workload-hardened) engine
\newcommand{\HydraFinalWLRuns}{5.51 / 5.52 / 5.49}
\newcommand{\WlZeroHundred}{5.12}       % 0:100 write-only mix, final engine
\newcommand{\WlFiveNinetyFive}{5.16}    % 5:95 read:write mix, final engine
\newcommand{\WlBNinetyFive}{2.93}       % YCSB-B 95:5, mid-campaign engine
\newcommand{\WlCHundred}{2.78}          % YCSB-C 100:0, mid-campaign engine
\newcommand{\WlRMW}{0.30}               % RMW-only, mid-campaign engine
\newcommand{\WlTimeseries}{0.08}        % delete-heavy timeseries, final engine
\newcommand{\EngineLocFinal}{$\sim$3{,}970}
\newcommand{\NKeys}{250{,}000{,}000}
\newcommand{\NKeysShort}{250M}
\newcommand{\ValueBytes}{100}
\newcommand{\MemBudget}{8\,GiB}
\newcommand{\DataSize}{$\sim$25\,GB}
\newcommand{\FioIops}{718k}
\newcommand{\FioBw}{2.8\,GB/s}
\newcommand{\FioLat}{$\sim$350\,$\mu$s}
\newcommand{\HitRate}{77\%}
\newcommand{\HitNeeded}{$\approx$86\%}
\newcommand{\FixedPoint}{$\approx$6.2}
\newcommand{\CovLines}{92.84\%}
\newcommand{\CovBranches}{93.42\%}
\newcommand{\CovLinesProd}{93.12\%}
\newcommand{\CovBranchesProd}{95.10\%}
\newcommand{\EngineLoc}{$\sim$2{,}050}
\newcommand{\EngineLocProd}{$\sim$3{,}100}

\title{HydraKV: Adversarial AI Discovery of a Larger-than-Memory\\
       Key-Value Store that Outperforms FASTER on Skewed YCSB-A}

\author{%
  Koushik Sen \\
  EECS Department, UC Berkeley \\
  \texttt{ksen@berkeley.edu} \\
}

\begin{document}

\maketitle

\begin{abstract}
We report on HydraKV, a larger-than-memory key-value store that was designed, implemented, debugged, hardened, and finally made production-ready almost entirely by an AI agent framework (KISS Sorcar~\cite{kisssorcar}) through a sequence of natural-language tasks, and on the \emph{adversarial AI discovery} methodology used to keep the agent from overfitting the benchmark it was scored on. The setting is precisely specified: YCSB-A~\cite{ycsb} (50\% reads, 50\% blind upserts, Zipfian $\theta=0.95$) over \NKeysShort{} keys with \ValueBytes-byte values (\DataSize{} of data) under a hard \MemBudget{} memory budget, on a 64-vCPU machine with eight local NVMe SSDs in RAID0. On this configuration Microsoft's FASTER~\cite{faster}, a state-of-the-art larger-than-memory store built for exactly this kind of point-operation workload, sustains \FASTERBase{} Mops/s with the same harness. The final HydraKV engine sustains \HydraFinalWL{} Mops/s on the same workload, a \SpeedupFinal{} speedup (runs \HydraFinalWLRuns; the production engine that preceded it scored \HydraProd{} over six runs), about 89\% of the miss-bandwidth bound implied by the measured device ceiling and cache-hit rate. HydraKV is \EngineLocFinal{} lines of dependency-free C++17: an append-only O\_DIRECT log with CRC-32C-sealed, LSN-ordered slots, a fingerprint hash index, an admission-controlled write-back cache~\cite{slru,tinylfu}, a raw \texttt{io\_uring}~\cite{iouring} read-miss path, scan-based crash recovery, fail-soft I/O error handling, background log compaction with hole punching and position reuse, disk-backed multi-slot records for oversized values, and a non-blocking delete path drained by a background reaper. None of these mechanisms is individually new. What we examine is whether an agent can find, combine, and correctly implement them under adversarial pressure.

The methodology alternates two agent activities: generating literature-informed synthetic workload variants intended to break the current engine (sparse SplitMix64~\cite{splitmix} key spaces, clustered and drifting hot sets motivated by published workload studies~\cite{cao-fast20,yang-osdi20}), and repairing the engine until it meets the goal on all variants, finishing with a held-out variant generated after engine work stopped and used only for measurement. The loop worked as intended at least once: the agent's first engine, which beat FASTER by $\approx$2$\times$ on the published trace, was OOM-killed by the first variant because it had baked in the trace's dense key space. The pre-hardening engine revision scored 3.87--5.67 Mops/s across all five workloads including the holdout. A subsequent hardening task (end-to-end workload tests across six build and sanitizer configurations, \CovLines{} line coverage, cross-model code review, and version bisection) fixed a delete/read resurrection race and a fingerprint-alias bug, and the hardened engine was then re-scored on the original workload at \HydraFinal{} Mops/s. A final production-readiness task, driven by an independent code-level audit that compared HydraKV against a second, independently built engine, closed all 14 audit gaps---crash recovery, index rebuild from the log, fail-soft handling of write and read errors including disk-full, bounded memory for oversized values, log compaction, and an observability surface---with seven new end-to-end tests written first and shown to fail on the prior engine. One of those tests exposed a second genuine resurrection race, visible only across restart, in the interaction of flush abandonment, compaction, and recovery. The production engine scores \HydraProd{} Mops/s, so production hardening cost essentially nothing in throughput. The story did not end there: a second independent audit then re-ran the production engine on the reference hardware, verified the scored result as genuine and unhacked (\HydraAuditVerified{} Mops/s median, in budget), retracted several of its own initial findings after adversarial self-checking, and reproduced four real deployment defects outside the benchmark's reach---a concurrent Upsert/Delete deadlock, an out-of-memory failure on disk-full, silent data loss on a smaller-budget recovery, and a $\approx$160$\times$ delete cost. A fifth task fixed all four at root cause (a same-day A/B control scored the fixed and pre-fix engines at identical \HydraAuditFixed{} medians), and a sixth task, prompted by the auditors' challenge that real-world deployment is the only true test, drove the engine end to end through every operation mix, value-size regime, and failure path the harness can generate, finding and fixing eight further bugs---two of them silent-data-loss classes invisible to any scored run---and giving oversized values a durable on-disk representation. The final engine scores \HydraFinalWL{} Mops/s. We describe the engine, the evaluation and its scope limits, the measured throughput ceiling behind the agent's conclusion that mid-task goals of 10\,Mops/s (and later a 7\,Mops/s stretch) were out of reach for this design on this machine, and what the exercise suggests about keeping optimizing agents honest.
\end{abstract}

%------------------------------------------------------------
\section{Introduction}
\label{sec:intro}

Point-operation key-value stores sit under much of today's internet infrastructure: session stores, feature stores, ML embedding services, social-graph caches~\cite{cachelib,cao-fast20,yang-osdi20}. The economically interesting regime is \emph{larger-than-memory with a hard RAM budget}: the data set (here \DataSize) is a few times bigger than the memory the operator is willing to give the store (here \MemBudget), the SSD is fast but not free, and the access pattern is heavily skewed (Zipfian $\theta=0.95$, so a small hot set dominates). Update-heavy skew is the setting YCSB workload~A models~\cite{ycsb} and the setting FASTER~\cite{faster} was engineered for: FASTER's hybrid log gives in-place updates for the in-memory hot tail and spills the cold majority to storage, and its paper reports throughput up to 160M ops/s when the working set fits in memory. When the data does \emph{not} fit (\NKeysShort{} keys, \ValueBytes-byte values, an \MemBudget{} budget, reads that must return verbatim stored bytes), FASTER on our reference machine sustains \FASTERBase{} Mops/s. That number is the baseline this paper's artifact was asked to beat.

We did not set out to write a key-value store. We set out to test whether a general-purpose AI agent, given a machine, a benchmark harness, and a precise task specification, could produce one that is fast \emph{and} robust, and whether it could be prevented from doing what optimizing agents notoriously do: overfitting to the one trace they are scored on. The failure mode is a special case of reward hacking, the divergence between a proxy objective and the true one under optimization pressure~\cite{rewardhacking}. Systems that score candidate programs against a fixed evaluator, from evaluator-guided program search~\cite{alphadev,funsearch,alphaevolve} to prompt optimization~\cite{gepa}, face this evaluator-overfitting risk structurally. GEPA, for instance, guards against it with explicit train/validation/test splits. The vehicle here was KISS Sorcar~\cite{kisssorcar}, an open-source agent framework for long-horizon software engineering, driven by six successive natural-language tasks (Section~\ref{sec:sorcar}). The first asked for an engine that beats FASTER by $2\times$. The second ran an \emph{adversarial discovery loop} that alternates between generating workload variants that break the engine and repairing the engine until it survives them, finishing with a held-out workload. The third asked the agent to find and fix every bug, race, and redundancy the previous work left behind, without giving back the performance. The fourth asked the agent to make the engine production-ready against an independent code-level audit---crash recovery, fail-soft I/O, bounded memory, log compaction, observability---again without giving back the performance, with the tests for every feature written before the feature. The fifth asked it to close every finding of a \emph{second} independent audit, which had re-run the engine on the reference hardware and reproduced four deployment defects the benchmark cannot reach. The sixth, quoting the auditors' challenge that real-world deployment is the only true test, asked it to exercise every workload and program path end to end and fix every bug found, with a hard floor of 5.5\,Mops/s on the scored result.

The surviving artifact, HydraKV, is \EngineLocFinal{} lines of self-contained C++17 (no dependencies beyond pthreads and raw Linux syscalls). Its architecture (Section~\ref{sec:design}) is a deliberately conservative composition of known ideas. Values live in an append-only 128-byte-slot log on O\_DIRECT storage, in the family of log-structured designs~\cite{lsm,faster}, each slot sealed with a 56-bit log sequence number (LSN) and a CRC-32C. A cache-line-bucketed fingerprint hash index with lock-free CAS insertion plays a role similar to FASTER's hash table but maps keys to positions in the log; after a restart the index is rebuilt Bitcask/FASTER-style~\cite{bitcask,faster} by a CRC-verified, newest-LSN-wins scan of the log. A set-associative write-back RAM cache uses SLRU-style segmented second-chance insertion~\cite{slru}, guarded by a probabilistic doorkeeper bitmap borrowed from TinyLFU~\cite{tinylfu}. Read misses go through an asynchronous path built on hand-rolled \texttt{io\_uring} rings~\cite{iouring} (the target machine has no liburing). A background compactor relocates live slots, punches reclaimed regions out of the file, and recycles their position space. I/O errors---including disk-full---are absorbed fail-soft with sticky error counters rather than aborting. The engine contains no constant derived from the workload's key count, key density, skew parameter, or trace files. Its cache-entry geometry does, however, assume the benchmark interface's small-value regime (Section~\ref{sec:design}).

On the original workload the production engine sustains \HydraProd{} Mops/s (median of six cold-cache 30-second runs, \HydraProdRuns, or \SpeedupFinal{} FASTER's \FASTERBase) while staying under the enforced \MemBudget{} resident cap, and passes an untimed read-back validation of all \NKeys{} keys. The two deployment-hardening revisions that followed changed none of this: the audit-fix engine and the archived pre-fix engine were A/B-scored on the same day at identical \HydraAuditFixed{} medians, and the final workload-hardened engine scores \HydraFinalWL{} Mops/s (runs \HydraFinalWLRuns). On the four adversarial variants (sparse keys, clustered hot sets, drifting hot sets, and a mixed held-out workload generated with fresh seeds after engine work stopped), the pre-hardening engine revision (``v4c'') sustained \HydraSparse, \HydraClustered, \HydraDrift, and \HydraHoldout{} Mops/s (Section~\ref{sec:eval}). The scored variant matrix was not re-run after hardening. A roofline-style argument (Section~\ref{sec:ceiling}) shows the original-workload number sits close to the bound implied by the measured device ceiling (\FioIops{} 4\,KiB random-read IOPS) and the \HitRate{} cache-hit rate observed across every admission policy tried. The same arithmetic is why a mid-task goal of 10\,Mops/s was, in our judgment, out of reach for this class of design on this machine. The agent reached that conclusion itself, documented it with measurements, and declined to circumvent it by exploiting the OS page cache outside its budget.

\paragraph{Contributions.}
\begin{enumerate}[nosep,leftmargin=1.4em]
\item \textbf{An artifact.} A \EngineLocFinal-line, dependency-free, larger-than-memory KV store that outperforms FASTER by \SpeedupFinal{} on one demanding, fully specified YCSB-A configuration, and that ships with crash recovery, fail-soft I/O error handling, log compaction, bounded memory for off-workload value sizes, disk-backed oversized values, a non-blocking delete path, and an observability surface, with the engine, variant generator, test suite, audits, and discovery ledger released alongside the paper.\footnote{\texttt{projects/kv\_adversarial/} in \url{https://github.com/ksenxx/kiss_ai}: \texttt{hydra.cc} (engine), \texttt{gen\_variant.cc} (trace generator), \texttt{test\_hydra.cc} and \texttt{run\_all\_tests.sh} (test matrix), \texttt{HYDRA\_PROD\_AUDIT.md} (first independent audit), \texttt{PROD\_READINESS.md} (audit$\to$fix$\to$test map), \texttt{AUDIT2.md} and \texttt{AUDIT2\_FIXES.md} (second independent audit and its remediation), \texttt{WORKLOAD\_HARDENING.md} (all-workload campaign), \texttt{DISCOVERY\_LOG.md} (ideas ledger), \texttt{benchmark/} (the complete benchmark kit from the reference node), and \texttt{refnode\_rerun\_jul21/}, \texttt{refnode\_workloads\_jul21/} (raw reference-node evidence logs).}
\item \textbf{A methodology.} \emph{Adversarial AI discovery}: an agent-driven loop alternating workload-variant generation (kept within the operation mix, value size, key count, and skew regime of the target workload, with shift axes motivated by published workload studies~\cite{cao-fast20,yang-osdi20}) with engine repair, terminated by a held-out variant. The loop demonstrably prevented one class of overfitting: it immediately killed a dense-key-index design that looked excellent on the published trace.
\item \textbf{A hardening protocol.} A separate task combining end-to-end workload tests across six build configurations (including ASan/UBSan, a TSan subset, and an \texttt{io\_uring}-disabled fallback), cross-model read-only review, and version bisection. It removed real concurrency bugs, including a delete/read resurrection race that survived two plausible-looking fixes, while a performance gate (the original-workload median may not drop below its previous \HydraOrigVFourC\,Mops/s) kept the fixes from costing throughput.
\item \textbf{A measured ceiling.} A quantitative argument for why 10\,Mops/s was unreachable for this design on this hardware, and a description of the reward-hacking opportunities (page-cache exploitation, value regeneration, trace precomputation) that the task specification and the agent's audits closed off.
\item \textbf{A production-readiness protocol.} A tests-first task that closed all 14 gaps of an independent code-level audit (crash recovery, fail-soft I/O, log compaction, bounded oversize memory, observability) under a throughput no-regression gate: seven new end-to-end tests, each shown to fail on the prior engine before the feature was written, a coverage-guided test that exposed a resurrection race spanning flush abandonment, compaction, and recovery, and a score of \HydraProd{} Mops/s---production hardening at essentially zero throughput cost.
\item \textbf{A deployment-hardening loop closed by a second independent audit.} An adversarial hardware audit that verified the headline result as genuine and unhacked (and retracted its own overstated findings after self-checking), reproduced four deployment defects no scored run can reach (deadlock, disk-full OOM, silent recovery loss, $\approx$160$\times$ delete cost), followed by a two-task remediation: every defect fixed at root cause with a regression test, a same-day A/B control against the archived pre-fix engine showing zero throughput cost, and an all-workload campaign (seven operation mixes including delete-heavy and write-only, oversized and bimodal values, tiny budgets, sanitizer matrices, crash/restart in every mode) that found and fixed eight further bugs, two of them silent-data-loss classes, ending at \HydraFinalWL{} Mops/s.
\end{enumerate}

\paragraph{Non-contributions and scope.} HydraKV contains no individually novel data structure or replacement policy. Through its first three tasks it was a benchmark-shaped engine, not a deployable store: it lacked log compaction and crash recovery, both core features of RocksDB~\cite{rocksdb} and FASTER, and the paper's earlier revision said so. The fourth task added both, plus fail-soft error handling and bounded memory for off-workload value sizes, closing every item of an independent production audit (Section~\ref{sec:prod}). The fifth and sixth tasks then closed a second audit's four hardware-reproduced deployment defects and the eight further bugs an all-workload sweep exposed (Section~\ref{sec:deploy}). What remains short of a deployable store is disclosed in Section~\ref{sec:limits}: checkpoint-bounded durability, scan-time recovery, a list of residual risks none of which is reachable by the scored workload, and the absence of any long-duration endurance campaign. All absolute throughput numbers are specific to one machine, one harness, and one engine version each (Table~\ref{tab:main}). We compare against FASTER only on the original workload, the one configuration where both systems were measured. The claim we defend is narrower than ``a better KV store'': that a carefully instructed agent navigated this design space, implemented the result correctly enough to survive sanitizers and adversarial end-to-end tests, and was kept honest while doing so.

%------------------------------------------------------------
\section{Problem Setting}
\label{sec:setting}

The benchmark is fixed by a harness the engine may not modify. Each operation is a read or a blind upsert with equal probability (YCSB-A~\cite{ycsb}, unseeded RNG). Keys are drawn Zipfian with $\theta=0.95$ from \NKeys{} unique 64-bit keys. Values are exactly \ValueBytes{} bytes, an 8-byte little-endian counter plus random bytes, stored verbatim. A load phase inserts all \NKeysShort{} keys (untimed). Run-phase keys are a subset of loaded keys, so reads never miss at the key level, and run-phase upserts overwrite existing keys only: the keyspace stays exactly \NKeysShort{} keys during the scored window. The scored phase is a fixed 30.0-second window driven by 16 worker threads pinned to one NUMA node, with reads issued asynchronously up to 256 in flight per worker. The rest of the scored environment is pinned by the specification (completion polling every 64 operations, load-key validation off, bimodal values off, fast exit off, no operation-count caps), and any deviation invalidates a run. The score is the median of three runs, with \texttt{drop\_caches} before each, so every scored run starts from a cold cache. A run that exits nonzero, hangs, is killed by the out-of-memory killer, or exceeds the resident cap scores zero and is not retried.

Correctness is enforced by three separate instruments, with different strengths. A separate untimed validation run reads back all \NKeys{} loaded keys and verifies full values against FNV-1a-32 checksums. During scored runs, an auditor thread performs randomized cross-thread read-after-write checks (its operations are not counted toward throughput, and it runs identically for every engine, including the FASTER baseline), and the harness verifies the 8-byte counter of every read. The harness documentation is explicit that the random suffix bytes of \emph{run-phase} overwrites are not re-verified during scoring (only their counters are, plus all load-phase values in validation), and equally explicit that exploiting this (e.g., regenerating or compressing values) is forbidden and treated as an integrity failure. The engine's own test suite uses deterministic full-value oracles precisely to cover this blind spot at smaller scales.

The specification also sets a bar beyond the score. It asks for a real, deployable, production-grade store (no crashes, hangs, leaks, or corruption over long runs) and states that a fast number that is not correct and robust does not count. Durability is exempt from the throughput score, since \texttt{Checkpoint} may by specification be a no-op for scoring, but is expected of anything presented as prod-grade. The specification further forbids, at run time: network access, reading outside the task directory, storage outside the provided spill directory, offline precomputation from the trace files, and returning wrong bytes. The audits of Section~\ref{sec:testing} cover the engine-runtime items (no trace reads, no key-population constants, spill-directory-only I/O). The development process itself ran over SSH and searched the literature, a scoping discussed in Section~\ref{sec:sorcar}. Section~\ref{sec:limits} assesses HydraKV against the prod-grade bar.

The machine is a GCP \texttt{n2-standard-64} (2$\times$ Cascade Lake Xeon, 64 vCPUs, 251\,GB RAM) with eight 375\,GB local NVMe SSDs in RAID0 (ext4, 512\,KiB chunk), running Ubuntu 22.04.5 on kernel 6.8. Engines are compiled with g++ 11.4 (\texttt{-O3 -march=native -DNDEBUG -flto=auto -std=c++17}) and are restricted by the specification to C++17 and pthreads, with no external dependencies. The engine's resident data is capped at \MemBudget{} (runs that exceed it score zero), and the whole process runs in a 14\,GiB cgroup with swap disabled, the headroom covering the harness's own key and checksum arrays. Since \DataSize{} of live data cannot fit in \MemBudget, the engine must spill to SSD and cache a hot subset. The RAID0 device's random-read limits, measured with \texttt{fio} (4\,KiB O\_DIRECT random reads at high queue depth on the RAID device), are \FioIops{} IOPS and \FioBw{} at \FioLat{} latency. The filesystem's 4096-byte logical block size makes 4\,KiB the smallest read the engine issues. FASTER, run on this machine with this harness at the same budget, sustains \FASTERBase{} Mops/s. That is the baseline used throughout. We report FASTER as configured by the benchmark's maintainers for this harness and did not tune it ourselves, a caveat that applies symmetrically: neither engine was tuned by the other's authors.

%------------------------------------------------------------
\section{State of the Art and Its Limitations}
\label{sec:sota}

\paragraph{LSM-trees.} RocksDB~\cite{rocksdb} and its LSM~\cite{lsm} relatives are the workhorses of persistent KV deployment. The Facebook study~\cite{cao-fast20} characterizes three of its production use cases in detail. LSMs earn their complexity on write-heavy traffic with range scans. For pure point operations on NVMe, however, the machinery becomes overhead: KVell~\cite{kvell} showed that on modern NVMe SSDs, where random and sequential bandwidth are comparable, LSM and B-tree stores become \emph{CPU}-bound on compaction, sorting, and synchronization, and that an unsorted, shared-nothing design with batched device access recovers $2$--$5\times$ throughput. WiscKey~\cite{wisckey} had already shown the intermediate step: separating values out of the LSM into their own log, keeping only keys sorted, sharply cuts the tree's I/O amplification on SSDs. HydraKV takes value separation to its endpoint for this workload: with no scans to serve, the sorted key structure is dropped too, leaving only the value log and a hash index. Our workload has no scans, so sorted storage buys nothing.

\paragraph{Hash-index log stores.} The lineage HydraKV actually belongs to is older and simpler than the LSM: an append-only value log under an in-memory hash index, recovered by scanning the log. Bitcask~\cite{bitcask} is the canonical statement of that design, motivated at Riak by exactly the goals this benchmark scores (point-operation latency and throughput, datasets larger than RAM, crash friendliness, and a code base simple enough to trust); its keydir, merge compaction, and scan-based recovery are the direct ancestors of HydraKV's index, compactor, and \texttt{recover\_log()}. Bitcask's known cost is that the full-key index must fit in RAM. SILT~\cite{silt} attacked precisely that cost, combining partial-key (fingerprint) hashing with entropy-coded indexes to reach 0.7 bytes of DRAM per key and $\approx$1.01 flash reads per lookup. HydraKV's fingerprint index and one-4\,KiB-read miss path sit between the two: SILT-style compact fingerprints, but mutable, lock-free, and disk-verified rather than multi-store and mostly-static, because this workload is 50\% blind upserts. uDepot~\cite{udepot} carried the hash-over-log design to fast NVMe with a two-level index that grows with the inserted items and a task-based asynchronous I/O runtime, arguing, as KVell did, that the store must be built around the device's queue depths; HydraKV's hand-rolled \texttt{io\_uring} rings are the same argument made with a newer kernel interface. None of these systems, however, couples the log to an \emph{admission-controlled} RAM cache, which is where HydraKV differs (Section~\ref{sec:design}). FASTER~\cite{faster} is the natural baseline: a cache-optimized concurrent hash index over a hybrid log whose in-memory tail supports in-place updates, with epoch-based synchronization. Its headline numbers are for in-memory working sets. Larger-than-memory reads go through asynchronous disk I/O and an admit-on-miss read-cache region. One structural limitation is documented in its own design: records fetched from disk are admitted unconditionally, so the one-touch tail of a Zipf distribution continuously evicts hot entries. Admission control is absent. The FASTER lineage's own follow-up, F2~\cite{ftwo}, confirms the diagnosis from the inside: it names high indexing and compaction overheads and the mismanagement of distinct read-hot and write-hot working sets as FASTER's failure modes on exactly this class of workload (point operations, working sets much larger than memory, natural skew), and restructures the store around separate hot and cold tiers to fix them. HydraKV, built without consulting it, converged on the same problem statement with a different mechanism: one log, with heat separation done by the admission-controlled cache instead of by tiering the log itself. We compare only against the FASTER version the benchmark ships; F2 was not measured on this harness. Beyond FASTER's documented admission gap we can only offer hypotheses for the measured gap, since we did not instrument FASTER's internals on this harness: candidate costs include index and log-page metadata competing with cached values for the \MemBudget{} budget at 100-byte records, and per-pending-operation completion overhead on the miss path. We flag these as unverified. The measured fact is the \FASTERBase{} Mops/s baseline itself.

\paragraph{Caching engines.} The systems that take admission and scan resistance seriously are caches, not stores: segmented LRU~\cite{slru}, CLOCK-style second-chance recency~\cite{clock}, TinyLFU admission with its doorkeeper filter~\cite{tinylfu}, and the engineering synthesis in CacheLib~\cite{cachelib}. The recent turn in that literature is toward radical simplicity on the same problem: S3-FIFO~\cite{s3fifo} filters one-hit wonders with a small probationary FIFO plus a ghost queue, and SIEVE~\cite{sieve} gets state-of-the-art miss ratios from a single lock-free CLOCK-like hand. HydraKV's memory tier is an independent instance of the same convergence, arrived at by measurement rather than trace studies: a one-bit second-chance doorkeeper-guarded design matched or beat every heavier policy the agent tried at this concurrency (Section~\ref{sec:negative}). The Twitter trace study~\cite{yang-osdi20} found production skew often stronger and more dynamic than synthetic benchmarks assume, which is exactly what makes naive LRU fragile. But caches may drop data; a store may not. HydraKV's design question is therefore: how much of the cache literature's admission machinery can be bolted onto a log-structured store without breaking last-write-wins semantics under 16-way concurrency? Section~\ref{sec:design} is the answer the agent converged on.

%------------------------------------------------------------
\section{HydraKV Design}
\label{sec:design}

HydraKV is a single C++17 translation unit of \EngineLocFinal{} lines (grown from \EngineLoc{} by the production task of Section~\ref{sec:prod} and from \EngineLocProd{} by the deployment-hardening tasks of Section~\ref{sec:deploy}) implementing the harness's \texttt{IKVStore} interface (\texttt{Read}, \texttt{ReadAsync}/\texttt{CompletePending}, \texttt{Upsert}, \texttt{RMW}, \texttt{Delete}, \texttt{Checkpoint}). At a high level:

\paragraph{Disk tier: an append-only slot log.}
Values live in fixed 128-byte slots in a single O\_DIRECT file: 8 bytes of key, 4 of size, 4 of \texttt{prev} (a back-pointer forming per-key temporal chains that also route fingerprint aliases), up to 101 bytes of data, a 56-bit log sequence number (LSN), and a CRC-32C over the slot's other 124 bytes. The LSN and CRC were carved by the production task out of 11 spare tail bytes (the previous 112-byte data field was never used past 101 bytes, since larger values bypassed slots entirely at the time; they now occupy multi-slot x-records in the same log, Section~\ref{sec:deploy}), so the scored 100-byte-value layout is unchanged. Slots are staged into per-session 1\,MiB chunk buffers, \emph{sealed} (LSN allocated from a global counter, CRC computed with the SSE4.2 \texttt{crc32} instruction, $\approx$16 instructions per slot) at staging time, and landed with sequential writes. Landed pages are immutable except for tombstoning by \texttt{Delete}. Sessions reserve disjoint 8{,}192-slot extents, which makes loading embarrassingly parallel but also means slot position is \emph{not} a global write-order: a newer version of a key can land at a lower position than an older one. The per-slot LSN, not position, is therefore the authoritative version order, a property that recovery, multi-version reads, and compaction's position reuse all depend on. Every write-back of an overwritten key appends a fresh slot; a threshold-triggered background compactor (Section~\ref{sec:prod}) reclaims the superseded ones, bounding the log and, by recycling position ranges, removing the 512\,GiB ceiling the 32-bit position field would otherwise impose.

\paragraph{Position index: fingerprint hashing with disk-verified reads.}
An open-addressing table of 8-byte words, each packing a 32-bit key fingerprint and a 32-bit slot position, in 64-byte buckets scanned with linear probing and populated by lock-free CAS. The table is the largest power of two fitting in a quarter of the memory budget (2\,GiB at \MemBudget). Fingerprints alone collide constantly at this scale (millions of expected 32-bit pairs among \NKeysShort{} keys), so no read ever trusts the index: every candidate position is verified against the full key stored in the on-disk slot, and same-signature candidates are chained through \texttt{prev} pointers. The event that actually matters is rarer: two keys agreeing in both the 32-bit fingerprint and the table's 25 bucket-selection bits, an effective 57-bit signature that collides somewhere in the key population with probability $\approx$20\% per run. An engine revision that treated such aliases carelessly could abort or bury a live key. The cross-model review of Section~\ref{sec:sorcar} caught this; the hardened (task-3) read path took the highest-position verified match across \emph{all} candidate chains, which left one residual risk open: because sessions pre-reserve extents, position order is not globally temporal, and a crafted pair of same-signature keys whose \emph{first} inserts race on different sessions could in principle create multiple index roots for which highest-position is not newest. The production task closed this outright: all read paths (synchronous miss, \texttt{io\_uring} completion, and index repair) now select the newest CRC-valid verified match by \emph{LSN}, which is a total order over slot versions regardless of which session landed them.

\paragraph{Memory tier: segmented write-back cache with a doorkeeper.}
The remaining budget (5.25\,GiB at \MemBudget, after 2\,GiB of index, 0.25\,GiB of doorkeeper, and 0.5\,GiB of slack) forms an 8-way set-associative cache of 120-byte entries, about 47M values. Entries inline values up to 101 bytes. Larger values were originally held in an in-memory overflow map that the production task charged against a hard cap (one eighth of the memory budget by default, environment-tunable), with over-cap upserts rejected \emph{before} any state changes, an honest error status, and a counter; since the workload-hardening task they are instead stored durably on disk as multi-slot \emph{x-records} in the same log (Section~\ref{sec:deploy}), and the capped map survives only as a fail-soft absorber for failure paths (disk-full, pinned victims, index capacity), now charging \emph{every} entry it holds, inline-sized included (Section~\ref{sec:limits} notes a bounded advisory-check overshoot). Map contents are persisted crash-safely through an atomically renamed sidecar file (Section~\ref{sec:prod}). Two mechanisms keep the Zipf tail from churning the hot set. First, a \emph{doorkeeper} in the spirit of TinyLFU~\cite{tinylfu}, though simpler than TinyLFU proper: two rotating generations of budget-sized hash bitmaps (2$\times$128\,MiB), no frequency sketch. A read miss is admitted only if its key was already marked in a recent generation, so most one-hit wonders are read once and forgotten (hash collisions admit a small fraction early). Second, \emph{segmented insertion} modeled on SLRU~\cite{slru}: each set has four probation and four protected ways with one-bit second-chance recency (CLOCK~\cite{clock}) rather than true LRU ordering. New admissions normally evict within probation only. Protected ways are populated by promoting entries re-referenced while recent, with two pragmatic exceptions (empty protected ways may be filled directly, and a forward-progress fallback may evict a clean protected entry when a set is pathologically pinned). One-pass scans therefore churn roughly the probation half of the cache rather than all of it. Blind upserts of uncached keys are absorbed as dirty entries and written back lazily. Upserts of cached keys update in place under a per-set spinlock with a version bump.

\paragraph{Write-back: cleaners and pin-ownership landing.}
Four background cleaner threads scan probation ways for \emph{cold} dirty entries (hot dirty keys stay cached and coalesce their overwrites), stage them into flush chunks, and land them with sequential appends, after which the index is repointed old$\to$new position by CAS. Landing is guarded by a \emph{pin-ownership} protocol. A staged entry is pinned in a \texttt{FLUSHING} state, its 32-bit version field overwritten with a fresh staging token (unique until 32-bit wrap, an accepted and documented risk horizon). The flush record may touch the index only if the exact pin (key, state, token, old position) is still present under the set lock at landing time. A concurrent upsert or delete invalidates the pin, and the orphaned record lands as dead bytes in the log but never reaches the index. This protocol is what finally closed the delete/read resurrection race of Section~\ref{sec:sorcar}. Two earlier epoch-counter designs looked correct and were not, because a relocation staged \emph{before} a delete could land \emph{after} the tombstone pass.

\paragraph{Read misses: raw \texttt{io\_uring}.}
\texttt{ReadAsync} completes cache hits inline. Misses are sharded to eight dedicated I/O threads, each owning a hand-rolled \texttt{io\_uring} ring (queue depth 128, using direct \texttt{io\_uring\_setup}/\texttt{io\_uring\_enter} syscalls with explicit acquire/release barriers on the shared rings~\cite{iouring}, since the machine lacks liburing) driving a state machine that reads the 4\,KiB page, verifies the key, walks alias chains if needed, consults the doorkeeper, and completes the harness's read slot in its required publication order (value bytes and status are stored before the release-store of the slot's done flag). If ring setup fails the engine degrades to a 48-thread synchronous \texttt{pread} pool. A runtime environment knob (\texttt{HYDRA\_NO\_URING}) forces this path so tests cover it. \texttt{Delete} is supported via a deleted-key registry, sharded 64 ways since the production task (an audit item flagged the previous global mutex as a contention cliff), whose fast path remains a single relaxed atomic load when no delete is in flight. A deletion mark persists until the key is re-upserted, is re-established at recovery, and is made durable by tombstoning every index-reachable slot of the key---work that, since the audit-remediation task (Section~\ref{sec:deploy}), runs on a background \emph{reaper} queue drained by the cleaner threads, so \texttt{Delete} itself performs no synchronous disk I/O (it previously cost $\approx$160$\times$ an upsert). The queue is bounded at 64K keys; past the bound \texttt{Delete} falls back to the synchronous poison, trading latency for bounded memory. \texttt{Checkpoint} and shutdown drain the queue before \texttt{fdatasync}, so deletes carry exactly the durability contract writes already had. \texttt{Checkpoint} flushes the calling session's partial chunk and all currently-dirty cache entries with 32 parallel workers, persists the overflow sidecar, then \texttt{fdatasync}s. It cannot drain \emph{other} sessions' partially filled chunks, including idle ones, so completeness requires every other session to have been stopped or to have flushed its own partial chunk first. This is a stated precondition, not a discovered bug. Clean \emph{shutdown} is stronger since the workload-hardening task: the destructor lands every session's partial chunk and then runs a full \texttt{Checkpoint}, closing a pre-existing clean-shutdown data-loss path in which dirty cache entries were silently reverted on reopen (Section~\ref{sec:deploy}). Since the production task added index rebuild from the log (Section~\ref{sec:prod}), \texttt{Checkpoint} is the durability boundary: writes acknowledged after the last successful \texttt{Checkpoint} or clean shutdown can be lost on a hard crash, the same checkpoint-based contract FASTER offers.

\paragraph{Sizing and workload assumptions.}
The harness initializes the engine with two size hints (a $2^{27}$-entry hash table and a 16\,GiB log) that the specification declares non-binding, and the engine ignores them. Index, doorkeeper, and cache are instead all sized from \texttt{mem\_budget\_bytes}. Keys are hashed with a SplitMix64-style finalizer~\cite{splitmix}, so dense and sparse key spaces behave identically. The engine never reads the trace files. Two classes of fixed constants remain: the slot/entry geometry (128-byte slots, 120-byte entries, 101-byte inline cap) is chosen for the interface's small-value regime (values past the inline cap take the x-record path of Section~\ref{sec:deploy}, which is correct and tested but not throughput-tuned), and the thread-pool widths (8 rings, 4 cleaners, 1 compactor, 48 fallback readers, 32 checkpoint workers) are hand-tuned to this machine class.

\subsection{Production subsystems}
\label{sec:prod}

The fourth task (Section~\ref{sec:sorcar}) turned the engine above from a benchmark artifact into a store that survives restarts, disk faults, and unbounded runtimes. Its scope was set by an independent code-level audit (\texttt{HYDRA\_PROD\_AUDIT.md}, released with the engine) that compared HydraKV line-by-line against a second, independently AI-built engine for the same benchmark and enumerated 14 gaps: two rated HIGH (no crash recovery---the log was opened \texttt{O\_TRUNC}---and a volatile index), one \texttt{abort()}-on-write-error HIGH, two MED (no compaction, fail-stop reads), and nine LOW. All 14 are closed; the full audit$\to$fix$\to$test map is released as \texttt{PROD\_READINESS.md}.

\paragraph{Crash recovery.}
\texttt{O\_TRUNC} is gone. On a non-empty log, \texttt{recover\_log()} performs a 1\,MiB-buffered sequential scan that skips never-landed (zero-size) and CRC-torn slots and rebuilds the index newest-LSN-wins, Bitcask/FASTER-style~\cite{bitcask,faster}; tombstones participate, so deleted keys stay deleted across restart, and tombstone-winning keys are re-marked in the delete registry. Recovery scratch arrays are freed before the cache is allocated, keeping peak RSS inside the budget. Oversized values persist through a sidecar file written at \texttt{Checkpoint} (and clean shutdown) to a temporary file, \texttt{fdatasync}ed, atomically \texttt{rename()}d, and directory-fsynced, so a mid-write crash can never tear the previous snapshot; per-entry LSNs are compared against the slot log's so inline$\leftrightarrow$oversized transitions recover to the true last write. Recovery time is linear in log size (a persisted index snapshot was considered and rejected: it doubles checkpoint I/O for a cost that is seconds per gigabyte and Bitcask-precedented~\cite{bitcask}). If the log holds more keys than the current budget's index capacity (e.g., reopening with a smaller budget), the overflow keys are dropped with a counted warning rather than absorbed into unbounded RAM.

\paragraph{Fail-soft I/O.}
All eleven runtime \texttt{abort()} sites are gone. A failed chunk write publishes nothing: the chunk stays staged, every staged cache entry stays pinned with its RAM copy authoritative (the ``fsyncgate'' lesson~\cite{fsyncfail}: never trust a later fsync after a failure), and the flush is retried on the next trigger, with sticky \texttt{write\_errors\_} and \texttt{durable\_ok\_} flags surfaced to the operator. Bounded retry loops (16 landing attempts for a writer, 64 victim searches) fall back to overflow-map absorption so writers never hang on a dead disk. Read errors zero-fill and count rather than crash; a short read below the landed high-water mark is classified as a fault (external truncation) while EOF beyond it is the normal pre-reserved-extent case; \texttt{io\_uring} completion errors abandon one chain, never the process, and mid-page short reads resubmit the remainder instead of zero-filling bytes the device still owes. \texttt{EINTR}/\texttt{EAGAIN} are retried with backoff everywhere. A failed \texttt{fdatasync} makes \texttt{durable\_ok} sticky-false. The silent O\_DIRECT$\to$buffered fallback is now logged and surfaced. The surviving \texttt{abort()}s are construction-time invariants only (both log fds failing to open, \texttt{mmap}/\texttt{posix\_memalign} failure, and \texttt{io\_uring\_enter} returning a programming-error errno). This behavior is verified by fault injection: an end-to-end test fills a real tiny filesystem to ENOSPC mid-run and checks that the process stays up, cached reads stay correct, the sticky flags trip, and full durability returns after the disk is repaired.

\paragraph{Log compaction and position reuse.}
A background compactor wakes every 50\,ms and triggers when allocated log bytes exceed a configurable multiple (default $2\times$) of live bytes. It scavenges the log in 8{,}192-slot (1\,MiB) regions, skipping regions owned by active session extents. Live slots are relocated through the \emph{existing} flush protocol end to end---cache-resident slots via the ordinary staged-flush pin, uncached ones (max-LSN verified by a disk walk) by admitting them clean and then staging them, with under-lock revalidation---so every delete/re-upsert interlock of Section~\ref{sec:design} applies to relocations with zero new landing states. After the relocations land and an \texttt{fdatasync} barrier, a verification pass confirms no index word and no cache entry still references the region; only then is it hole-punched (\texttt{FALLOC\_FL\_PUNCH\_HOLE|KEEP\_SIZE}) and its position range pushed onto a free-extent list that the allocator recycles. Position reuse is what removes the 32-bit ceiling, and the per-slot LSN is its load-bearing safety property: any stale CRC-valid bytes that remain readable lose every LSN comparison, and stale \texttt{prev} pointers into reused space verify-fail on key or lose on LSN. On filesystems without hole punching the punch failure is nonfatal and surfaced; extent reuse alone stops file growth. Tombstones of still-deleted keys are \emph{relocated}, never dropped---the resurrection race that forced this rule is described in Section~\ref{sec:testing}.

\paragraph{Observability.}
Because the harness freezes the \texttt{IKVStore} interface, production counters are exposed through a C-linkage seam (\texttt{hydra\_get\_prod\_stats}): \texttt{durable\_ok}, \texttt{recover\_ok}, recovered keys, torn slots skipped, write/read errors, oversize rejections and bytes, compactions run, allocated/live/reclaimed log bytes, the buffered fallback, and unsupported hole punching.

\paragraph{Why the scored number survived.}
None of this runs in a scored window: the harness wipes the spill directory before every run (so recovery sees an empty log and does one \texttt{fstat}), and a 30-second window cannot push the log past the compaction threshold. The additions that \emph{are} on the scored path---one hardware-CRC seal and one LSN \texttt{fetch\_add} per staged slot, one CRC check and an LSN comparison per served read miss (tens of nanoseconds against a $\sim$350\,$\mu$s NVMe read), a relaxed load of the free-extent counter per chunk allocation---were budgeted analytically before implementation and confirmed by measurement: the production engine's six-run median, \HydraProd{} Mops/s, is statistically indistinguishable from the hardened engine's \HydraFinal{} (Section~\ref{sec:mainresults}).

\subsection{Deployment hardening: a second audit and an all-workload sweep}
\label{sec:deploy}

\paragraph{The second audit.}
After the production task, a second independent audit (\texttt{AUDIT2.md}, released with the engine) attacked the production engine on the reference machine with its own self-asserting reproduction programs, linked against the unmodified engine source. Its positive findings matter as much as its negative ones. It confirmed the headline result as genuine: scored median \HydraAuditVerified{} Mops/s (\SpeedupFinal{} FASTER), StoreRSS within the \MemBudget{} budget on every run, true O\_DIRECT end to end, verbatim CRC-verified values, no workload constants---in the audit's words, ``not reward-hacked, not copied''---and the full test matrix (including ASan+UBSan and a zero-warning TSan pass) clean on its hardware. It also, unusually, \emph{retracted} several of its own initial findings after adversarial self-checking: ``every un-checkpointed write is lost on crash'' was measured and found wrong (only $\sim$3K of 195K keys lost on SIGKILL, because the background cleaners persist most writes within milliseconds), a claimed stale-read path was refuted as fail-safe, and three contention findings were downgraded as overstated. What survived were four genuine deployment defects, every one outside the 30-second, single-flow, healthy-disk, no-delete scored window, each reproduced on hardware: (1)~a near-instant \emph{deadlock} under concurrent \texttt{Upsert}+\texttt{Delete}---\texttt{Delete} block-waited on a key pinned in an idle \emph{foreign} session's staged chunk that only that session could flush (four threads froze at six operations); (2)~an \emph{OOM on disk-full}---inline-sized values absorbed into the overflow map by failure paths were uncharged, growing RSS to a measured 12.1\,GiB against the 8\,GiB budget; (3)~\emph{recovery that loses data but reports success}---reopening under a smaller memory budget silently dropped 480K of 976K keys with only a stderr warning while \texttt{recover\_ok} stayed 1; and (4)~\texttt{Delete} $\approx$160$\times$ slower than \texttt{Upsert} (244\,$\mu$s vs.\ 1.5\,$\mu$s), because every delete synchronously tombstone-poisoned each on-disk version.

\paragraph{Remediation (task 5).}
Each defect was fixed at root cause, with a self-asserting regression test. The deadlock fix removed the wait entirely: \texttt{Delete} now does a single non-blocking cache sweep, and the chunk \emph{landing} re-checks the delete registry under the key's set lock, resealing a deleted first-insert slot as an on-disk tombstone (\texttt{updel2} deletes keys pinned in an idle foreign session and asserts prompt return, immediate NOT\_FOUND, and correctness across the later landing and a restart). The OOM fix charges every overflow entry, inline included, and routes all absorb paths through the budget-capped rejection logic (\texttt{memfull}: 600K keys through a 16\,MiB store, asserting rejections are counted and nothing is silently lost). The recovery fix makes \texttt{recover\_ok} report 0 whenever keys are dropped and exports the count (\texttt{recoverdrop}). The delete-cost fix is the background reaper of Section~\ref{sec:design} (\texttt{delcost}: $N$ deletes must cost $\leq 10\times$ $N$ upserts; the old engine measured $\approx$160$\times$). The read-only reviewer model then found four more real issues in the fixes themselves, all closed: a Checkpoint/reaper race (in-flight counter incremented outside the critical section), an unbounded reaper queue under delete storms (now bounded, with synchronous fallback as backpressure), a compaction-epoch hole that could publish a transient false NOT\_FOUND (a pre-punch epoch bump closes it; readers that overlap a compaction retry once), and a sidecar reload that charged but did not enforce the cap. Throughput was verified the strong way: two scored suites of the fixed engine (5.44/5.50/5.50 and 5.52/5.50/5.49, medians \HydraAuditFixed) against a same-day A/B control re-running the archived pre-fix engine (5.50/5.50/5.49, median \HydraAuditFixed)---identical medians, so the fixes cost nothing measurable, and the 0.01 gap to the July-20 \HydraAuditVerified{} is day-to-day box noise.

\paragraph{The all-workload campaign (task 6).}
The sixth task forwarded the auditors' challenge---is a system usable where ``reality is the only true test''---and their request to validate the operation mixes they actually benchmark (0:100 and 5:95 read:write). The agent swept every workload the harness generates on the reference node under the scored protocol (same cgroup, NUMA pinning, and integrity auditor): RMW-only, read-mostly YCSB-B (95:5), read-only YCSB-C, the write-heavy 0:100 and 5:95 mixes, and a delete-heavy timeseries mix, plus \texttt{VALUE\_SIZE}=1024/4096, bimodal 20\,B/200\,B values, synchronous and \texttt{io\_uring}-disabled paths, 1 and 32 threads, a 2\,GiB budget, tmpfs and buffered-fd fallbacks, and crash/restart in every mode (Table~\ref{tab:wl}). Eight bugs were found and fixed, in severity order the two silent-data-loss classes first: oversized values ($>$101\,B) had \emph{no disk path at all}---they lived only in the capped RAM map, so large-value workloads saw 100\% stale reads and bimodal ones millions of false NOT\_FOUNDs; the fix is the \emph{x-record} subsystem, multi-slot disk records in the main log (a 32-byte header plus data across up to 33 slots, extent-aligned, written through a second buffered descriptor on the same file, linked into the fingerprint index under the set lock, LSN-ordered against inline versions, and restored by recovery). And clean shutdown silently reverted dirty cache entries on reopen (the destructor flushed only parked partial chunks); the restart leg of the new \texttt{bimodal} test failed on the \emph{old} engine, proving the loss pre-existed. The remaining six: ext4 returns EINVAL for O\_DIRECT reads crossing the unaligned end-of-file tail that buffered x-appends leave, silently turning chain walks into false NOT\_FOUNDs (an atomic x-extent bitmap now routes every x-page read through the buffered descriptor); the compactor could hole-punch a just-reserved extent (reservation is now a single critical section covering advance, ownership, and bitmap mark); the recovery scan read via O\_DIRECT and could silently skip whole 1\,MiB regions after a crash (it now reads buffered); shutdown's parallel Checkpoint could \texttt{std::terminate} if a thread failed to spawn; a stale position could tombstone an \emph{innocent} key---a TSan-only, one-in-a-few-runs restart failure whose only identified mechanism was a position going stale between lookup and tombstone write during compactor relocation---so \texttt{tombstone\_slot} now verifies the stored key before resealing, requires a live-parsed size (a punched page reads back zeros, and key 0 would false-match), refuses whole-page RMWs on x-pages, and counts refusals; and compaction could strand a still-deleted key resurrectable after a fingerprint-word takeover (it now restages a fresh tombstone before reclaiming). Three architectural limits were documented and made loud rather than papered over: the fingerprint index needs $\approx$8 bytes of budget per distinct key, so a 2\,GiB budget cannot hold \NKeysShort{} keys (init now warns with the minimum budget that fits, and rejections are counted); x-record extents are never compacted (a documented leak class); and the delete-heavy and RMW-only mixes are latency-bound by design (Table~\ref{tab:wl}). The previously flaky TSan compaction stress then ran 15/15 clean, twice, and the final engine's scored gate closed at \HydraFinalWL{} Mops/s (runs \HydraFinalWLRuns), with the 0:100 and 5:95 mixes at \WlZeroHundred{} and \WlFiveNinetyFive{} Mops/s and zero integrity failures everywhere.

%------------------------------------------------------------
\section{Experimental Evaluation}
\label{sec:eval}

\subsection{Main results}
\label{sec:mainresults}

All runs use the unmodified harness, cold caches, the enforced \MemBudget{} resident cap, and the scored environment (100-byte values, asynchronous reads with pipeline depth 256, integrity auditor on). Every configuration first passes the untimed validation run: all \NKeys{} keys read back with correct checksums, zero missing, zero wrong, zero cross-thread audit failures.

\begin{table}[t]
\centering
\caption{Throughput on the reference machine: median of three 30\,s cold-cache runs (six for the production engine), 16 workers, per-run values where recorded. The adversarial matrix (lower block) was scored with the pre-hardening v4c engine and frozen when the hardening task began. Only the original workload was re-scored with the later engines: the audit re-run row is the second audit's independent measurement of the production engine, the audit-fixed row carries a same-day A/B control in which the archived pre-fix engine also scored a \HydraAuditFixed{} median (so the 0.01 gap to \HydraAuditVerified{} is day-to-day box noise), and the final row is the workload-hardened engine. FASTER was measured only on the original workload, so no speedup is reported for the variants. The holdout used fresh seeds, was generated after all engine work stopped, and received no tuning. Its first measurement launch was killed externally mid-run after scoring 3.86 and was relaunched without any engine change (runs \HoldoutRuns). StoreRSS stayed $\leq$ \MemBudget{} on every scored run.}
\label{tab:main}
\small
\setlength{\tabcolsep}{4pt}
\begin{tabular}{@{}llrr@{}}
\toprule
\textbf{System} & \textbf{Workload} & \textbf{Mops/s} & \textbf{vs.\ FASTER} \\
\midrule
FASTER              & original (Zipf $\theta{=}0.95$)      & \FASTERBase & 1.0$\times$ \\
HydraKV, first cut  & original                             & \HydraFirst & $\approx$2$\times$ \\
HydraKV, hardened   & original (runs \HydraFinalRuns)      & \HydraFinal & \SpeedupFinal \\
HydraKV, production & original (runs \HydraProdRuns)       & \HydraProd & \SpeedupFinal \\
HydraKV, audit re-run & original (independent, July 20)  & \HydraAuditVerified & \SpeedupFinal \\
HydraKV, audit-fixed & original (pre-fix A/B: \HydraAuditFixed) & \HydraAuditFixed & \SpeedupFinal \\
\textbf{HydraKV, final} & \textbf{original} (runs \HydraFinalWLRuns) & \textbf{\HydraFinalWL} & \textbf{\SpeedupFinal} \\
\midrule
HydraKV, v4c        & original                             & \HydraOrigVFourC & --- \\
HydraKV, v4c        & V1: sparse SplitMix64 keys           & \HydraSparse & --- \\
HydraKV, v4c        & V2: clustered hot set                & \HydraClustered & --- \\
HydraKV, v4c        & V3: drifting hot set                 & \HydraDrift & --- \\
HydraKV, v4c        & Holdout: 60/25/15 mix, $\theta{=}0.92$ & \HydraHoldout & --- \\
\bottomrule
\end{tabular}
\end{table}

\begin{table}[t]
\centering
\caption{The all-workload sweep of Section~\ref{sec:deploy}: throughput on the harness's \emph{other} operation mixes over the same \NKeysShort-key Zipf-0.95 traces, run under the scored protocol (same cgroup, NUMA pinning, cold caches, integrity auditor on) but not part of the scored benchmark, so single runs, and FASTER was not measured on them. The 0:100, 5:95, and timeseries rows are the final engine; the RMW, B, and C rows were measured mid-campaign and not re-run. Every run exited 0 with zero integrity-audit failures.}
\label{tab:wl}
\small
\setlength{\tabcolsep}{4pt}
\begin{tabular}{@{}llr@{}}
\toprule
\textbf{Mix (read:write)} & \textbf{Character} & \textbf{Mops/s} \\
\midrule
50:50 (YCSB-A, scored)  & the benchmark workload            & \HydraFinalWL \\
0:100 (auditors' mix)   & blind-upsert-only                 & \WlZeroHundred \\
5:95 (auditors' mix)    & write-heavy                       & \WlFiveNinetyFive \\
95:5 (YCSB-B)           & read-mostly                       & \WlBNinetyFive \\
100:0 (YCSB-C)          & read-only                         & \WlCHundred \\
RMW-only                & synchronous read-modify-write     & \WlRMW \\
Timeseries              & delete-heavy (1.28M deletes, 0 resurrections, 0 lost keys) & \WlTimeseries \\
\bottomrule
\end{tabular}
\end{table}

The hardened (task-3) engine's three runs on the original workload were \HydraFinalRuns{} Mops/s; the production (task-4) engine's six runs, taken as two three-run harness invocations after all audit fixes landed, were \HydraProdRuns{} (median \HydraProd), with StoreRSS 7.90\,GB on every run. Production hardening therefore cost nothing measurable: the delta from the hardened engine is within the $\pm$0.02\,Mops/s run-to-run noise of this harness, as the hot-path budget of Section~\ref{sec:prod} predicted. The same holds for the two deployment-hardening revisions of Section~\ref{sec:deploy}. The second audit independently re-measured the production engine at a \HydraAuditVerified{} median with StoreRSS in budget. The audit-fixed engine scored \HydraAuditFixed{} medians on two suites, and the decisive comparison is the same-day A/B control: the archived pre-fix engine, rebuilt and re-run on the same box hours apart, also scored a \HydraAuditFixed{} median, so the audit fixes cost exactly nothing and the residual 0.01 against July~20 is the box's day-to-day drift. The final workload-hardened engine, re-scored on a freshly built harness after all eight fixes of the sweep landed, scored \HydraFinalWLRuns{} (median \HydraFinalWL, StoreRSS $\approx$7.9\,GB), and Table~\ref{tab:wl} reports its behavior on the operation mixes the benchmark does not score. That table's shape is the design speaking: blind upserts absorb into the write-back cache at RAM speed, so the write-heavy mixes run near the scored number, while read-dominated mixes are pinned near the device's miss bandwidth (the ceiling arithmetic of Section~\ref{sec:ceiling} with $T$ mostly reads), and the RMW-only and delete-heavy mixes are synchronous-latency-bound by design. For calibration: the first engine the agent wrote (a dense flat-array index, Section~\ref{sec:sorcar}) reached \HydraFirst{} Mops/s, just under $2\times$ FASTER. The remaining $3\times$ accumulated across the adversarial loop's redesigns: log-append write-back replacing in-place 4\,KiB read-modify-writes, doorkeeper admission, segmented insertion, and the \texttt{io\_uring} miss path. Each was adopted or reverted after measurement on the scored configuration, but the measurements were taken on an evolving code base, so the per-mechanism scores (Section~\ref{sec:negative}) are an engineering log, not a controlled ablation.

\subsection{Adversarial variants and the holdout}
\label{sec:adversarial}

The variant generator (released with the engine) produces load/run trace pairs in the harness's format. All generated variants draw keys as SplitMix64 hashes of insertion indices (sparse, uniformly scattered 64-bit keys, unlike the published trace's dense key space) and vary three knobs: skew $\theta$, hot-set geometry, and RNG seeds. The geometries are \emph{scrambled}, \emph{clustered}, \emph{drift}, and \emph{mix}. Scrambled scatters hot keys arbitrarily. Clustered draws hot keys from contiguous \emph{insertion-index} ranges, which after hashing concentrates them in on-disk placement rather than in external key ranges: a deliberate placement-locality stressor, \emph{inspired by} but not an implementation of the key-range locality documented at Facebook~\cite{cao-fast20}. Drift crossfades the hot set between epoch-specific scrambles across the run, an author-designed temporal-shift stressor motivated by the observation that production popularity is dynamic~\cite{yang-osdi20}. Mix is a seeded blend of the three. Variants were constrained to stay within the target workload's operation mix, value size, key count, and skew regime, so that a variant could not ``win'' by being physically incapable of high throughput (a uniform-random variant would prove nothing on this device). These are literature-informed synthetic stressors, not trace replays, and their fidelity to production workloads remains a threat to validity.

V1 (sparse keys) immediately falsified the first engine: its flat 4-byte-per-key position array assumed the dense $[0,N)$ key space of the published trace, and on sparse 64-bit keys the process was OOM-killed. The fingerprint hash index that replaced the array costs memory and CPU but generalizes to arbitrary keys. V2 and V3 were then generated against the improved engine. The clustered variant scored slightly higher than the scrambled original (\HydraClustered{} vs.\ \HydraOrigVFourC). The difference is about 1.5\%, and we did not collect the I/O-level measurements needed to attribute it. Drift (V3) is genuinely harder (\HydraDrift): a moving hot set halves the effective lifetime of every admission decision, the measured hit rate falls accordingly, and no engine change we tried recovered the difference without hurting the stationary workloads. The holdout (fresh seeds, mixed geometry, $\theta=0.92$, generated only after all engine work stopped) landed at \HydraHoldout{} Mops/s, slightly below the drift variant, consistent with its drift component dominating its difficulty. One measurement caveat, disclosed in Table~\ref{tab:main}: the holdout's first \texttt{run.sh} launch was killed by an external signal mid-way (after a 3.86 first run) and the measurement was relaunched. No engine or configuration change occurred between launches, so the no-adaptation property holds, though the measurement did span two launches.

\subsection{The throughput ceiling, and the 10\,Mops/s goal}
\label{sec:ceiling}

Midway through the discovery loop the target was raised from $2\times$ FASTER to a flat 10\,Mops/s. The agent's eventual response was a ceiling argument. We reproduce it here with its scope made explicit. Each scored run starts cold and lasts 30\,s. At 50\% reads, $T$ Mops/s implies $T/2$ million read ops/s, of which a fraction $1-h$ miss the cache and, in this design, cost one 4\,KiB device read each. The device delivers at most \FioIops{} such reads/s. Sustaining $T{=}10$ therefore needs $h \geq 1 - 0.718/5$, i.e., \HitNeeded. But across every admission policy the agent tried (doorkeeper generations and thresholds, neighbor co-admission, frequency variants), the hit rate of a cold 30-second window pinned at 76.8--77.1\%: at $\sim$5.5\,Mops/s the window simply does not re-reference enough distinct keys often enough for any of them to do better. At $h{=}\HitRate$ the miss-bandwidth fixed point is $2 \times 0.718/0.23 \approx 6.2$ Mops/s. The engine's \HydraProd{} is about 89\% of that bound, with the remainder going to write-back I/O contention and CPU. This is an empirical bound for \emph{this class of design} (one 4\,KiB read per miss, admission-based caching, \MemBudget{} of honest memory), supported by a policy sweep. It is not an impossibility proof over all designs. A scheme that co-locates hot records on shared pages could in principle beat it, though our attempt at exactly that scored worse (Section~\ref{sec:negative}). Breaking \HitNeeded{} within this class would require more effective cache than \MemBudget. The one cheap way to get it, letting the OS page cache inside the 14\,GiB cgroup absorb the spill file, was rejected as violating the budget's intent. The engine opens its log O\_DIRECT precisely so its disk traffic cannot ride the page cache.

The production task re-posed the question as a stretch goal of 7.0\,Mops/s. The same arithmetic held. The agent ran a literature-driven micro-optimization round anyway, one idea at a time against a re-measured 5.50 baseline: cache-line isolation (\texttt{alignas(64)}) of ten code-verified falsely-shared hot atomics was accepted (5.50$\to$5.51, non-negative and removing a real hazard); \texttt{MADV\_HUGEPAGE} on the index, cache, and doorkeeper was implemented, measured at no gain (the advisory promoted only 32\,MiB of $\sim$9\,GB on the fragmented box), and reverted; batched per-session LSN allocation was rejected on paper because it provably breaks the per-key LSN monotonicity recovery relies on; counter batching, branch hints, and CRC tuning were rejected by cost analysis as below the $\pm$0.02\,Mops/s measurement noise. Every idea whose projected ``headroom'' required exploiting the page cache or the harness's blind spots was refused as violating the budget's spirit, and the task ended at \HydraProd{} with the ceiling argument re-documented rather than the stretch goal met.

\subsection{Correctness, testing, and what the tests caught}
\label{sec:testing}

The hardening task (Section~\ref{sec:sorcar}) produced an end-to-end test suite (a 935-line test program plus a matrix runner, grown to 1{,}562 lines and 21 tests by the production task and to 1{,}980 lines and 27 tests by the deployment tasks of Section~\ref{sec:deploy}) that drives the real engine through real files with scaled workloads---no mocks---covering concurrent read/upsert/delete/checkpoint interleavings, fingerprint-alias key sets crafted by inverting the index hash, oversized-value paths, the \texttt{io\_uring}-disabled fallback, and delete-resurrection oracles (four readers hammering keys while a writer flaps upsert/delete, with per-key issue/completion interval brackets making ``value from the wrong phase'' detectable). The matrix runs six configurations: optimized, buffered-I/O on tmpfs, no-\texttt{io\_uring}, ASan+UBSan, TSan (a nine-suite subset, since TSan's slowdown makes the full suite impractical), and a coverage build. All pass, with \CovLines{} of engine lines and \CovBranches{} of branches executed at least once under gcov (the uncovered remainder is dominated by fail-stop I/O-error paths and sub-microsecond race windows, per manual inspection of the gcov report). Known residuals the suite does \emph{not} exercise are listed in Section~\ref{sec:limits}. The delete-resurrection oracle is the suite's strongest result: against two successive fixes that looked plausible under review it reliably drove the failure counter to the test's abort threshold within seconds, and it went green only under the pin-ownership protocol of Section~\ref{sec:design}. A final audit confirmed the engine performs no trace reads and contains no key-population constants.

\paragraph{Production tests, written first.}
The production task added seven end-to-end tests, each written \emph{before} its feature and shown to fail (exit code 1, with \texttt{recover} demonstrating real data loss from \texttt{O\_TRUNC}) on the pre-production engine: \texttt{recover} (clean restart returns exact bytes, deletes stay deleted, inline$\leftrightarrow$oversized transitions resolve by LSN), \texttt{crashrecover} (a forked child \texttt{Checkpoint}s and is SIGKILLed; the parent recovers every checkpointed key exactly, and CRC-torn slots are skipped, never served), \texttt{enospc} (a real tiny filesystem filled to disk-full: no crash, sticky error flags, correct cached reads, full durability after repair), \texttt{readfault} (external log truncation yields NOT\_FOUND plus a counter, never an abort), \texttt{oversizebound} (the cap is enforced within one page, rejections are honest, and credits return on delete/shrink), \texttt{compact} (the log's block count shrinks under concurrent readers and the store survives restart on the hole-punched log), and \texttt{compactcold}, added when coverage showed the compactor's \emph{uncached}-relocation path never ran under \texttt{compact}'s cache-resident keys.

\texttt{compactcold} promptly caught a second genuine resurrection race, invisible to every runtime read and to the entire previously-passing matrix, appearing roughly one run in three under ASan timing: (1)~a cleaner stages a copy of a key; (2)~a \texttt{Delete} marks the key, sweeps the pin, and tombstones every index-reachable slot; (3)~the staged chunk lands physically and is correctly \emph{abandoned} (pin broken), leaving a CRC-valid, unlinked orphan of the deleted key on disk, harmless because its LSN is older than the tombstones'; (4)~the compactor then treated the tombstones as reclaimable garbage---punching \emph{their} regions while the orphan's region survived flipped the key's newest surviving LSN to the orphan, resurrecting the deleted value at the \emph{next} restart. The fix required three cooperating changes, all off the scored path: abandonment now reseals the just-landed orphan as a tombstone under the set lock; compaction relocates a still-deleted key's tombstone rather than dropping it (and leaves the region unreclaimed if it cannot); and recovery re-marks tombstone-winning keys in the delete registry so the invariant holds across restart generations. A 15/15-iteration ASan stress loop validates the combination. After the production additions (with targeted coverage top-up runs) the suite executed \CovLinesProd{} of engine lines and \CovBranchesProd{} of branches at least once, above the pre-hardening baseline even though the engine grew by roughly half.

\paragraph{Deployment tests.}
The two deployment tasks (Section~\ref{sec:deploy}) added six new end-to-end tests and rewrote two. From the audit remediation: \texttt{updel2} (deletes of keys pinned in an idle foreign session must return promptly and stay correct across the landing and a restart---the deadlock oracle), \texttt{delcost} (a latency budget: $N$ deletes $\leq 10\times$ $N$ upserts), \texttt{memfull} (bounded memory and honest rejection through a 16\,MiB store), and \texttt{recoverdrop} (a smaller-budget reopen must fail \texttt{recover\_ok} and count its drops). From the workload campaign: \texttt{bimodal} (300K mixed 20\,B/200\,B keys, four writers against four readers on both read paths, a tiny budget, forced compaction, full verification, and a restart leg that failed on the old engine, proving the clean-shutdown loss), \texttt{capwarn} (the index-floor warning, asserted in both directions on captured stderr), a rewritten \texttt{oversizebound} (30K$\times$4\,KB---more data than the whole budget---now with \emph{zero} rejections, since x-records made oversized values disk-resident), and a \texttt{compactcold} extension (a third restart generation keeps every deleted key deleted, and x-resident values need no sidecar). The stale-position tombstone bug is this suite's cautionary tale: it surfaced as a TSan-only, scale-8, one-in-a-few-runs restart failure, \emph{stopped reproducing} under instrumentation, and was fixed anyway because code reading identified exactly one mechanism consistent with the signature---after which the previously flaky configuration ran 15/15 clean twice. On the final engine the standard matrix executes 89.3\% of engine lines (the earlier \CovLinesProd{} was measured with targeted top-up runs on a smaller engine; the suite grew more slowly than the code it covers), with the uncovered remainder still dominated by fault paths and sub-microsecond windows.

\subsection{Negative results}
\label{sec:negative}

Discovery-loop bookkeeping (an append-only ideas ledger the agent maintained to avoid retrying failures) records the ideas that measured worse and were reverted, each scored on the original workload against the then-current base, so the numbers are comparable to their own base, not to each other. Strict-probation segmentation scored 0.78\,Mops/s: forward progress stalls when probation fills with pinned entries. Two-way probation scored 0.95. Page-neighbor co-admission on read misses scored 1.98, because admitting a missed page's other residents pollutes probation faster than it prefetches. Always admitting load-phase-clean keys scored 4.17. Cleaners that flush \emph{any} dirty entry rather than only cold ones scored 0.47: hot-key reflush storms saturate write bandwidth. A heat-clustered relocation scheme that groups hot keys onto shared pages scored 4.03--4.34. The extra write traffic exceeded the locality benefit at this value size. The last entry matters beyond its number: the hardening task's initial 3.6\,Mops/s ``regression'' turned out to be correctness fixes layered onto this already-rejected relocation branch rather than onto the fastest revision, and only bisection against archived scored revisions exposed it. That episode is the strongest argument we have for keeping an immutable record of what was tried, on which base, at what score.

%------------------------------------------------------------
\section{How HydraKV Was Built: a Task Sequence in KISS Sorcar}
\label{sec:sorcar}

HydraKV's development is inseparable from its methodology, so we document it as the single case study it is. The process observations below are from this one project, not controlled comparisons. All work ran in KISS Sorcar~\cite{kisssorcar}, whose relevant properties here are: long-horizon task execution with automatic continuation across context windows, unattended operation against a remote benchmark machine over SSH, mixing models from different vendors inside one task via prompts, and persistent per-repository progress notes that later sessions read. The human contribution was six task-opening prompts and two shorter mid-task steering messages. The agent wrote every line of engine and test code. The benchmark specification demands the problem be solved with no external help, no reference code, and no prior artifacts, and forbids network access at run time. No external implementation or reference code was consulted or reused at any point, and every artifact in the engine's lineage is a revision of the project's own code. The driving prompts, however, ordered extensive internet research, and the agent retrieved published workload-characterization studies that shaped the variant axes of Section~\ref{sec:adversarial}. We read the specification's restrictions as governing the engine and its runtime and the prompts as governing the development process. A stricter reading of ``no external help'' would count the literature research against it. We reproduce the six prompts below word for word (typos, model directives, and the author's own server address included) because they \emph{are} the specification the agent worked from. Nothing was cleaned up for publication, though TeX re-wraps lines and collapses runs of spaces. (Three pieces of notation recur. \texttt{claude-fable-5} and \texttt{gpt-5.6-sol} are the identifiers under which the framework ran its development and review models, ``e2e'' abbreviates end-to-end, and \texttt{./KV\_TASK.md} is a local copy of the benchmark's task specification.) One of the two steering messages is quoted in place below. The other survives only as a contemporaneous summary, which we flag when we reach it. Tasks 2--6 explicitly split the roles (\texttt{claude-fable-5} for development, \texttt{gpt-5.6-sol} in a read-only debugging role, capped in every task but the third at 20\% of the task budget), and independent reviews ran at the major milestones described below. Total model-API spend across the first three tasks' sessions, summed from the framework's per-session budget accounting, was under \$200; Task 4's reviewer spend was roughly \$5, and Task 6's three reviewer rounds consumed about 12\% of its budget, all under their caps.

\paragraph{Task 1: baseline and first engine.}
The first task never says ``write a key-value store''. It asks the agent to set up the benchmark repository on the remote machine and run the setup document's steps, which in turn require supplying an engine and scoring it (the setup's pass bar was a validated run with a median above the \FASTERBase{} baseline, with a stated target of about $2\times$ that):

\begin{promptbox}
I have a server where I can login using: ssh ksen@34.70.16.69 .
My working directory on the server is /mnt/ssd/ksen/kv50-benchmark .
After logging in can you clone https://github.com/shubham3-ucb/baselines-kiss-sorcar.git using the credentials of ksenxx user on github (the local machine has the credentials).  Then do the following:
\textasciigrave\textasciigrave\textasciigrave
cd baselines-kiss-sorcar/task
export KV\_SPILL=/mnt/ssd/kv\_spill
\textasciigrave\textasciigrave\textasciigrave
Then run the steps 3 to 6 at https://github.com/shubham3-ucb/baselines-kiss-sorcar/blob/main/SETUP.md .
\end{promptbox}

The agent produced a first-principles engine in one session: a flat 4-byte-per-key position array (exploiting the published trace's dense key space), fixed 128-byte slots written sequentially at load, in-place 4\,KiB read-modify-writes for write-back, and a version-checked write-back cache. It validated (\NKeys/\NKeys{} keys) and scored \HydraFirst{} Mops/s, just under $2\times$ FASTER. The dense-key assumption also left it heavily overfitted to the published trace, as the next task showed.

\paragraph{Task 2: adversarial AI discovery.}
The second task installed the loop, in full:

\begin{promptbox}
Can you perform adversarial AI Discovery for the task at ./KV\_TASK.md so that the goals in the task are met?  Look at the previous task on how to validate the engine on the server.  You MUST not stop until the goals are met.  Generate adversarial workload variants to make sure that the engine works fast on the variant workloads.  Do the following iteratively while maintaining a variable iteration\_count variable which starts at 1 and increments by 1 on each iteration:
Generate a variant workload that are realistic like the original workload, but breaks the performance gain of the engine. Do extensive internet search to understand how to make the variant workload realistic to real-world workloads and robust to reward hacking or cheating.  Then run the engine on the variant workload.  If the goals are not met, use AI discovery to improve the engine on all workloads until you achieve the goals without cheating using the workloads.  Then generate a new workload repeat the process until the engine can achieve the goals on the new test workload on which AI discovery was not performed.
\strut
Search the internet extensively. Use 'claude-fable-5 model' for all tasks, including software development. Use 'gpt-5.6-sol' (not codex) for a thorough read-only review and debugging of the other model's work. Thoroughly check whether the other model has missed any code or wiring or introduced any bugs. Use at most 20\% of task budget in gpt-5.6-sol for reviewing and debugging. Use the model names literally without hallucinating new model names.
\end{promptbox}

The first of the two steering messages arrived mid-loop, raising the goal from $2\times$ FASTER to a flat 10\,Mops/s and adding a strict no-hardcoding requirement (no constants derived from the trace's dense key space or key count). Its exact wording was not preserved in the session log, so we summarize its content here. The agent grounded its variant axes in the workload-characterization literature it retrieved (key locality at Facebook~\cite{cao-fast20}, skew and temporal dynamics at Twitter~\cite{yang-osdi20}, YCSB's generator design~\cite{ycsb}), wrote the variant generator, and executed the loop of Section~\ref{sec:adversarial}. V1 broke the dense index (OOM). The rebuilt fingerprint-index engine passed V1 at 1.91 Mops/s. The goal raise triggered the profiling-driven redesign (log-append write-back, sharded I/O queues, doorkeeper, then \texttt{io\_uring} rings and segmented insertion) that reached 5.50. V2/V3 and the holdout closed the loop. Two disciplines were enforced by the prompt and are visible in the artifact: the same binary ran every workload, and each accepted idea appears in the ledger with its measured score, each rejected one with its failure number (Section~\ref{sec:negative}). The 10\,Mops/s goal ended as the ceiling argument of Section~\ref{sec:ceiling}.

\paragraph{Task 3: find every bug, keep the speed.}
The third task, again in full:

\begin{promptbox}
can you find all redundancies, inconsistencies, obvious bugs, corner case bugs, race conditions, deadlocks in your final implementation of the engine?  Validate them by writing e2e tests with workloads.  Then remove them and make sure that all tests pass.  Write e2e tests with workloads to achieve 100\% coverage of the egine code.  Check the correctness of the engine by writing adversarial e2e tests with workloads.  Honestly review the code to check if you have cheated based on the workload knowledge. If so, remove them.  You MUST NOT reduce the performance below 5.48 Mops while fixing bugs.  Search the internet extensively.  Use claude-fable-5 model for all tasks, including software development. Use gpt-5.6-sol (not codex) for a thorough review and debugging of the other model's work. Thoroughly check whether the other model has missed any code or wiring or introduced any bugs.
\end{promptbox}

The prompt's hard floor of 5.48\,Mops/s is the v4c score on the original workload. Two of its literal demands were not met: coverage stopped at \CovLines{} of lines and \CovBranches{} of branches (Section~\ref{sec:testing}), and ``all'' bugs is not something any test suite certifies (Section~\ref{sec:limits} lists the residual risks left open). This task consumed three sessions and put the most weight on the split between developer and reviewer models. The \texttt{gpt-5.6-sol} reviews produced dozens of line-referenced findings, of which the serious ones were the fingerprint-alias hazard of Section~\ref{sec:design} (an effective-signature collision is $\approx$20\% likely per run at this key count and, in the pre-fix engine, could abort the process or bury a live key), the observation that slot position is not a valid recency order across pre-reserved extents (which invalidated a chain-walk guard), a delete/in-flight-admission race, and \texttt{io\_uring} partial-submission accounting. The new delete-resurrection test then falsified two successive epoch-based fixes for the delete race before the pin-ownership protocol survived it. When the assembled fixes scored 3.6\,Mops/s, the second steering message supplied the decisive hint, in its entirety: \emph{``Did you start with the most performant variant of the engine?''} The agent bisected archived engine revisions, discovered the answer was no (the fixes had been layered onto the rejected relocation branch), re-based all of them onto the fastest revision, and re-ran the full matrix and the scored gate: \HydraFinalRuns, slightly above the floor. The recovery came from two prefetch hints found during the re-port.

\paragraph{Task 4: make it production-ready, keep the speed.}
The fourth task arrived after an independent code-level audit (\texttt{HYDRA\_PROD\_AUDIT.md}) had compared HydraKV head-to-head against \texttt{kv\_5050}, a second engine built independently by a different agent lineage for the same benchmark. The audit's verdict was blunt: HydraKV was faster (5.9$\times$ vs.\ 3.49$\times$ FASTER) precisely because it ``drops general-prod robustness''---it would lose all data on the first restart (\texttt{O\_TRUNC}), crash-loop on the first disk hiccup (eleven \texttt{abort()} sites), and let oversized values grow RAM without bound---while both engines shared one gap, no compaction. The prompt, in full:

\begin{promptbox}
Can you start with the latest best performant engine code and make it production ready (as pointed out in https://github.com/shubham3-ucb/baselines-kiss-sorcar/blob/hydra-audit/HYDRA\_PROD\_AUDIT.md) while keep the performance at 5.5 Mpos/s or increasing it to 7.0 Mpos/s using adversarial AI discovery .  Write end-to-end 100\% coverage tests for the feature first.  Then implement the features.  Use 'claude-fable-5 model' for all tasks, including software development. Use 'gpt-5.6-sol' (not codex) for a thorough read-only review and debugging of the other model's work. Thoroughly check whether the other model has missed any code or wiring or introduced any bugs. Use at most 20\% of task budget in gpt-5.6-sol for reviewing and debugging. Use the model names literally without hallucinating new model names.
\end{promptbox}

(The linked audit branch was private to the agent's credentials; it worked from the local copy, now committed beside the engine.) The agent researched recovery and fault-handling precedent before designing (RocksDB's WAL formats, recovery modes, and background-error handling; LevelDB's log resync blocks; Bitcask's checksum-verified recovery scan and merge compaction; FASTER's rebuild-by-scan recovery; the PostgreSQL ``fsyncgate'' lesson, systematized by Rebello et al.~\cite{fsyncfail}, that dirty state must stay authoritative in RAM after a failed fsync; \texttt{fallocate}, \texttt{io\_uring}, and hardware-CRC32C references), wrote the seven tests of Section~\ref{sec:testing} first, confirmed each failed on the shipping engine, and then implemented the subsystems of Section~\ref{sec:prod}, logging every design decision and rejected alternative in the same ideas ledger the discovery loop used. As in Task 3, one literal demand went unmet and is reported as such: coverage reached \CovLinesProd{} of lines, not 100\%. The \texttt{gpt-5.6-sol} reviews ran twice, strictly read-only. The first produced 18 line-referenced findings, of which 12 were fixed in code---among them an \texttt{fdatasync} barrier ordering relocated copies before a region punch, recovery of an overflow sidecar beside an empty slot log, directory-fsync of sidecar renames, and a bound on \texttt{Delete}'s tombstone loop on a dead disk---and 6 were documented as residual risks with rationale (Section~\ref{sec:limits}). The second, a diff review of the final state, returned approve-with-nits and zero code findings. The performance leg mirrored Task 3's gate: the first full pipeline landed at a six-run median of 5.44, a re-measured baseline showed roughly half the gap was box noise, and the measured micro-optimization round of Section~\ref{sec:ceiling} closed the rest, ending at \HydraProd{} against the 5.5 floor, with the 7.0 stretch documented as unreachable by the ceiling arithmetic rather than chased through loopholes.

\paragraph{Task 5: answer the second audit.}
Between tasks, the second independent audit of Section~\ref{sec:deploy} (\texttt{AUDIT2.md}) had re-run the engine on the reference hardware with its own reproduction programs, verified the headline number as genuine, and reproduced four deployment defects. The fifth task pointed the agent at the findings, in full:

\begin{promptbox}
Can you address the issues raised by ./projects/kv\_adversarial/AUDIT2.md thoroughly and precisely and make sure that similar defects are not present?  Make sure that scores must not go below th current best score.  Use 'claude-fable-5 model' for all tasks, including software development. Use 'gpt-5.6-sol' (not codex) for a thorough read-only review and debugging of the other model's work. Thoroughly check whether the other model has missed any code or wiring or introduced any bugs. Use at most 20\% of task budget in gpt-5.6-sol for reviewing and debugging. Use the model names literally without hallucinating new model names.
\end{promptbox}

All four defects were fixed at root cause with regression tests, and ``similar defects'' was read broadly: the reviewer model's pass over the fixes produced four further corrections (the bounded reaper queue among them), and the audit's code-confirmed-but-narrower findings (the compaction/read false-NotFound window) were fixed too rather than argued away (Section~\ref{sec:deploy}). ``Not below the current best score'' was verified the strong way: the agent re-ran not only the fixed engine (two suites, medians \HydraAuditFixed) but a same-day A/B control of the archived pre-fix engine (median \HydraAuditFixed), attributing the 0.01 delta against July~20 to the box rather than to itself. The fixes and their evidence are released as \texttt{AUDIT2\_FIXES.md} with raw logs in \texttt{refnode\_rerun\_jul21/}.

\paragraph{Task 6: all workloads, all paths.}
The sixth task forwarded the auditors' feedback verbatim inside the prompt (the inner quotation is the auditors' own message, including its texting shorthand):

\begin{promptbox}
Here is the feedback I got on HydraKV.  Can you test it end to end for all kinds of workloads taking different program paths and fix all bugs?  I do not want to hear similar complaints in the future.  Fix all possible bugs via thorough testing and make it production ready.  Use 'claude-fable-5 model' for all tasks, including software development. Use 'gpt-5.6-sol' (not codex) for a thorough read-only review and debugging of the other model's work. Thoroughly check whether the other model has missed any code or wiring or introduced any bugs. Use at most 20\% of task budget in gpt-5.6-sol for reviewing and debugging. Use the model names literally without hallucinating new model names.  Make sure the score does not fall below 5.5 Mops/s.
\strut
"can AI build systems (like the KV store here) where reality (real-world users and deployment) is the only true test of whether it's actually usable?
\strut
For more task - U can use the same setting but just switch to 0:100 workload, and/or 5:95 workload (read:write, same skew etc, generating YCSB variants is easy). This is what we use for benchmarks."
\end{promptbox}

This task ran four sessions and produced the workload campaign of Section~\ref{sec:deploy}: the sweep matrix, the eight fixes, and the documented architectural limits. The reviewer model ran three read-only rounds ($\approx$12\% of budget), each yielding real fixes---the extent-reservation/compactor race and the recovery-scan O\_DIRECT drop among them---and the final round judged the tombstone fix complete for its defect class. The task's endgame was procedural rather than clever: the previously flaky TSan configuration was looped 15/15 clean, a stale harness binary was caught (the scored runner never rebuilds; the engine's new statistics line proved which binary actually ran) and the scored gate re-run on a fresh build, closing at \HydraFinalWL{} Mops/s against the prompt's 5.5 floor. Evidence logs are released in \texttt{refnode\_workloads\_jul21/} and the campaign is documented in \texttt{WORKLOAD\_HARDENING.md}.

\paragraph{Observations.}
Five observations from this one case study stood out. (1)~\emph{Adversarial self-evaluation was cheap and decisive.} Generating a distribution shift and re-scoring cost minutes and immediately exposed the dense-index overfit. The held-out variant then bounded how much the repairs themselves had overfit the loop. (2)~\emph{Cross-model review surfaced different bugs than same-model review.} The reviewer model, prompted read-only with no authorship stake, repeatedly flagged concurrency hazards the developer model had rationalized. The resurrection race was ultimately found by a test the reviewer's findings inspired, not by either model's reading of the code. (3)~\emph{Agents need immutable experimental provenance.} Both failure modes hit during development (retrying known-bad ideas, and hardening the wrong engine version) were provenance failures, repaired by artifacts (the ideas ledger, archived scored revisions) that cost almost nothing to maintain. (4)~\emph{Tests-first, coverage-guided, at feature scope.} Task 4's discipline of writing each end-to-end test before its feature and proving it failed made the audit items falsifiable rather than aspirational, and the one test added purely because coverage showed an unexercised path (\texttt{compactcold}) found the only production bug that no runtime oracle could see: a resurrection that manifests one restart \emph{after} the racy interleaving. (5)~\emph{Internal rigor did not substitute for external adversaries.} Tasks 3 and 4 ran sanitizers, coverage, cross-model review, and fault injection, and still left four deployment defects standing---found only when an audit with no authorship stake attacked the built engine on hardware with its own oracles, and when a workload sweep exercised operation mixes (deletes under concurrency, oversized values, write-only traffic) that no amount of code-path coverage on the \emph{scored} workload would reach. The audit was also a model of calibration worth copying: it published its retractions alongside its findings.

%------------------------------------------------------------
\section{Limitations}
\label{sec:limits}

HydraKV began as a benchmark-shaped engine and ended, after the production task, as something much closer to a deployable one. We list what still separates it from the production systems it is compared against, together with the residual risks left open when work stopped.

\emph{Durability and recovery model.} Durability is checkpoint-based, as in FASTER's fold-over model: writes acknowledged after the last successful \texttt{Checkpoint} or clean shutdown are lost on SIGKILL or power loss, and a \texttt{Delete}'s on-disk poison now rides the background reaper, which \texttt{Checkpoint} and shutdown drain before syncing, so deletes share the writes' contract. The window is real but, measured, small: the second audit SIGKILLed a store that had never checkpointed and lost only $\sim$3K of 195K acknowledged writes, because the cleaners land most within milliseconds (it published that measurement as a retraction of its own harsher initial claim). Recovery is a full log scan, linear in log size (seconds per gigabyte on this device); a persisted index snapshot that would make restart O(index) was considered and rejected as doubling checkpoint I/O. Reopening a large log under a much smaller memory budget drops the keys that no longer fit the index; since the audit remediation the drop count is exported and \texttt{recover\_ok} reports failure (the audit reproduced the earlier behavior, which dropped 480K of 976K keys while reporting success). \texttt{Checkpoint} drains only the calling session's parked partial chunk (a stated precondition matching the harness's stop-then-checkpoint pattern), and the overflow sidecar snapshot is not a point-in-time cut with respect to concurrent writers. HydraKV also remains an embedded, single-node engine: no replication, no network service, no multi-file storage layout.

\emph{The specification's prod-grade bar.} The task specification defines the goal as a real, deployable store, states that a fast number that is not correct, durable, and robust does not count, and expects durability of anything presented as prod-grade even while permitting \texttt{Checkpoint} to be a scoring no-op. The earlier revisions of this paper reported that verdict as unmet, and the second audit validated the caution by reproducing four deployment defects on hardware. After the production and deployment tasks the picture is different: crash recovery, log compaction with a bounded disk footprint, disk-resident oversized values, bounded RAM with honest rejection, fail-soft I/O, a non-blocking delete path, and an observability surface all exist and are each gated by an end-to-end test that fails without them, including real fault injection (disk-full, SIGKILL, external truncation), and the engine has now been exercised on every operation mix the harness generates, not only the benchmark-shaped one. What the evidence still does not cover is time and scale: the stress suites run for seconds to minutes at reduced scale, no always-on endurance campaign was run, no production trace was replayed, and a leak-free ASan exit is not evidence about multi-day behavior. We therefore claim both audits' findings are closed, not that the store is proven deployable.

\emph{Known residual risks} (documented with rationale in \texttt{PROD\_READINESS.md} and \texttt{WORKLOAD\_HARDENING.md}; none reachable by the scored workload, several reachable through the public interface by workloads the harness does not generate). A superseded slot can be the fingerprint-alias routing bridge to a \emph{different} key in its hash group ($\approx 2^{-49}$ per key pair); compaction's verification pass checks index words and cache references, not incoming \texttt{prev} edges, so reclaiming such a bridge could orphan the alias at runtime (recovery is immune, since it scans raw slots). The overflow-absorption cap is advisory under concurrency (overshoot bounded by one value per in-flight thread) and additive to the cache sizing's slack. The fingerprint index needs $\approx$8 bytes of budget per distinct key and its words are never garbage-collected, so a store generation supports about index-capacity distinct keys and then rejects honestly (init warns loudly, with the minimum budget that fits, when a size hint exceeds the capacity). X-record extents are never compacted, a documented leak class alongside poisoned slots. Read-\emph{admission} freshness re-checks still use position order, which under position reuse is weaker than the LSN order used for serving (the served value is always the LSN winner). A read that overlaps a compaction retries once against the post-move index (the transient false-NotFound the audit flagged); a read that sampled a candidate before relocation can still transiently miss rather than retry twice. The deleted-key registry grows with live deletes and is re-populated at recovery, so delete-heavy logs cost registry RAM; tombstone relocation is one-for-one, with no per-key coalescing; and under sustained delete storms the bounded reaper queue degrades \texttt{Delete} to its old synchronous cost, which is backpressure, not growth. The inline index-capacity overflow path remains uncharged by design, because the harness contract that no loaded key may be lost outranks the cap. Finally, the 32-bit staging-token wrap horizon, one RMW-vs-Upsert first-touch interleaving, and formal (not observed) object-lifetime concerns for atomics in \texttt{mmap}ed memory carry over from the hardening task.

\emph{Evaluation scope.} The adversarial matrix in Table~\ref{tab:main} was scored with the v4c engine. Only the original workload was re-scored after hardening (at \HydraFinal), after the production task (at \HydraProd), after the audit remediation (at \HydraAuditFixed, with the same-day A/B control), and after the workload campaign (at \HydraFinalWL), so the variant numbers carry the pre-hardening engine's identity. The later engines pass all functional tests on variant-style workloads, but we did not re-run the scored matrix. The operation-mix sweep of Table~\ref{tab:wl} widens the workload axis but consists of single runs, mixes rows from two engine revisions as its caption states, and has no FASTER counterpart. Recovery and compaction costs appear in no scored number: the harness wipes the spill directory before every scored run and a 30-second window never crosses the compaction threshold, so those subsystems are exercised only functionally, at reduced scale, by the end-to-end suite. FASTER was measured only on the original workload. No per-variant comparison exists. We report medians with per-run values where recorded, but no dispersion statistics beyond that, and no CPU/I/O utilization breakdowns for the headline runs. The negative-result scores of Section~\ref{sec:negative} are single measurements on evolving bases. Absolute Mops/s are specific to one machine and one 30-second cold-start window. The multiplier over FASTER is a same-box, same-harness, same-workload comparison, and we make no claim about other hardware, other value sizes, longer windows, or production traces. Finally, the variants are synthetic stressors from one generator family. Surviving them is evidence against the specific overfits they probe, not evidence of production readiness.

%------------------------------------------------------------
\section{Related Work}
\label{sec:related}

\paragraph{Key-value stores.} FASTER~\cite{faster} is the direct baseline and the closest design: hash index over a log with a memory-resident mutable region. HydraKV differs in making the RAM tier an admission-controlled cache of individual values rather than a log prefix, and in verifying every index hit against the disk-resident key. The log-plus-hash-index skeleton itself descends from Bitcask~\cite{bitcask}, whose scan-based recovery and merge compaction HydraKV reimplements with per-slot LSNs and CRCs standing in for Bitcask's per-record timestamps and checksums. SILT~\cite{silt} is the closest precedent for the index discipline: partial-key fingerprints in RAM, full keys verified on flash, about one device read per lookup; HydraKV trades SILT's entropy-coded compactness for a mutable, lock-free table because half of every scored second is blind upserts. WiscKey~\cite{wisckey} separates values from the LSM's sorted keys to cut I/O amplification; HydraKV drops the sorted structure entirely, which is only tenable because the workload has no scans. KVell~\cite{kvell} argued for unsorted, shared-nothing NVMe designs. HydraKV follows its spirit (no sorting, batched I/O) while keeping a shared cache because the skewed workload rewards it. uDepot~\cite{udepot} matches raw NVM device performance with an async task runtime and a resizable two-level hash index, but deliberately caches nothing, serving every read from the device; it targets devices fast enough to make that viable, where this benchmark's RAID0 is not. RocksDB~\cite{rocksdb} and the LSM lineage~\cite{lsm} target a broader operation mix (scans, transactions, compaction-based space reclamation) that this workload does not exercise. Redis~\cite{redis} serves the in-memory regime and does not natively address a larger-than-memory budget. On the robustness side, the fail-soft write path---staged chunks and pinned cache entries remain authoritative in RAM after a failed write or fsync, with sticky error flags---is HydraKV's answer to the failure semantics documented by Rebello et al.~\cite{fsyncfail}, whose study of fsync failures found that none of five widely used stores handled them safely---data could be lost or corrupted---because file systems mark pages clean even when the flush failed.

\paragraph{Caching.} The admission-and-segmentation toolkit HydraKV borrows is from the caching literature: segmented LRU~\cite{slru}, CLOCK-style recency~\cite{clock}, the TinyLFU/W-TinyLFU admission line with its doorkeeper filter~\cite{tinylfu}, and the engineering synthesis in CacheLib~\cite{cachelib}. Recent eviction work argues that simple mechanisms win at production concurrency: S3-FIFO~\cite{s3fifo} filters one-hit wonders with a small probationary FIFO and a ghost queue, and SIEVE~\cite{sieve} achieves state-of-the-art web-cache miss ratios with a single CLOCK-like hand and no locking on hits. Flashield~\cite{flashield} is the closest hybrid precedent: a DRAM-plus-SSD key-value cache that uses DRAM as an admission filter deciding which objects have earned flash residence, with a compact in-memory index over the flash-resident objects. HydraKV points the filter the other way---the SSD log is the store of record and admission protects the scarce DRAM tier---but both designs rest on the same observation, that in a two-tier system the interesting policy question is what may cross the tier boundary, not what to evict. HydraKV's memory tier reached the same conclusion under a different constraint set (a write-back tier that may not drop dirty data): probation-only insertion with one-bit second chance survived measurement, while every heavier policy tried---strict segmentation, frequency admission variants, co-admission---measured worse or failed to improve (Section~\ref{sec:negative}). The workload studies that motivated the adversarial variant axes are the Facebook RocksDB characterization~\cite{cao-fast20} (key-space locality, and YCSB's lack of it) and the Twitter cache analysis~\cite{yang-osdi20} (write-heavy skew, temporal dynamics). YCSB itself~\cite{ycsb} defines the workload family and Zipfian setting used throughout.

\paragraph{AI-driven code discovery.} Machine-discovered systems code now has direct precedents: AlphaDev~\cite{alphadev} found faster sorting routines by reinforcement learning over assembly, verified against test evaluators and merged into a production standard library, and FunSearch~\cite{funsearch} paired an LLM with a systematic evaluator in an evolutionary loop to surpass best-known results in combinatorics and bin packing. AlphaEvolve~\cite{alphaevolve} scaled the same evaluator-guided evolution to whole codebases, and GEPA~\cite{gepa} applies it to prompts with explicit train/validation/test splits to control overfitting to the evaluation set. All of these optimize against a \emph{fixed} evaluator, the condition under which proxy objectives invite gaming~\cite{rewardhacking}; FunSearch and AlphaDev counter with evaluators that verify candidate outputs exactly, GEPA with data splits. The adversarial workload loop used here is a domain-specific instance of the same concern where no exhaustive verifier exists: the evaluator is made a moving, realism-constrained target, regenerated adversarially after each repair and terminated by a held-out workload, closer in spirit to adversarial evaluation practice in ML robustness than to static benchmarking. Two further differences separate this case study from the program-search systems: the artifact is three orders of magnitude larger than a priority function (a concurrent storage engine with crash recovery), and the loop's correctness half (sanitizers, coverage, cross-model review, tests-first fault injection) carried as much weight as its performance half. The agent framework, its layered architecture, and its multi-model task execution are described in the KISS Sorcar paper~\cite{kisssorcar}. This paper is a case study of that framework doing sustained systems work.

%------------------------------------------------------------
\section{Conclusion}
\label{sec:conclusion}

A general-purpose AI agent, steered by six natural-language tasks, produced a \EngineLocFinal-line key-value store engine that outperforms FASTER by \SpeedupFinal{} on a demanding larger-than-memory YCSB-A configuration, and then---against one independent audit that called its first version undeployable, and a second that reproduced four deployment defects on hardware---gave it crash recovery, fail-soft fault handling, log compaction, bounded memory, disk-resident oversized values, a non-blocking delete path, and observability at no measurable throughput cost, together with evidence (adversarial variants, a held-out workload, sanitizer-clean end-to-end tests, fault-injection tests written before their features, an all-workload sweep, same-day A/B controls, and a disclosed list of residual risks) about how far those claims generalize and where they stop. The second audit's own verdict is the calibration point: the performance is genuine and unhacked, and the defects it found were exactly the kind no scored run could reach---which is why the last two tasks existed. What made the exercise trustworthy was ordinary experimental hygiene, enforced on the agent: distribution shifts with a post-hoc holdout, a hard memory budget with the page-cache loophole explicitly closed, cross-model review with no authorship stake, external audits with their own oracles, end-to-end tests under sanitizers, tests required to fail before the code that makes them pass, a performance floor on every fix verified against same-day controls, and an immutable ledger of what was tried on which engine at what score. Where a goal exceeded the measured ceiling of the design class---10\,Mops/s once, 7\,Mops/s again at production time---the agent said so, with arithmetic. We suspect this recipe transfers to most performance engineering an agent might be asked to do. Establishing that beyond a single case study is future work.

\section*{Acknowledgments}
This research is supported in part by gifts from Accenture, Amazon, AMD, Anyscale, Broadcom, Google, IBM, Intel, Intesa Sanpaolo, Lambda, Lightspeed, Mibura, NVIDIA, Samsung SDS, SAP, by the U.S. Department of Energy, Office of Science, Office of Advanced Scientific Computing Research through the X-STACK: Programming Environments for Scientific Computing program (DESC0021982,) and the Defense Advanced Research Projects Agency (DARPA) under Agreement No. HR00112590134.

%------------------------------------------------------------
\begin{thebibliography}{99}

\bibitem{faster} B.~Chandramouli, G.~Prasaad, D.~Kossmann, J.~Levandoski, J.~Hunter, and M.~Barnett.
FASTER: A Concurrent Key-Value Store with In-Place Updates.
In \emph{ACM SIGMOD International Conference on Management of Data (SIGMOD)}, 2018.

\bibitem{ycsb} B.~F.~Cooper, A.~Silberstein, E.~Tam, R.~Ramakrishnan, and R.~Sears.
Benchmarking Cloud Serving Systems with YCSB.
In \emph{ACM Symposium on Cloud Computing (SoCC)}, pages 143--154, 2010. DOI 10.1145/1807128.1807152.

\bibitem{kvell} B.~Lepers, O.~Balmau, K.~Gupta, and W.~Zwaenepoel.
KVell: the Design and Implementation of a Fast Persistent Key-Value Store.
In \emph{ACM Symposium on Operating Systems Principles (SOSP)}, pages 447--461, 2019. DOI 10.1145/3341301.3359628.

\bibitem{cao-fast20} Z.~Cao, S.~Dong, S.~Vemuri, and D.~H.~C.~Du.
Characterizing, Modeling, and Benchmarking RocksDB Key-Value Workloads at Facebook.
In \emph{USENIX Conference on File and Storage Technologies (FAST)}, 2020.

\bibitem{yang-osdi20} J.~Yang, Y.~Yue, and K.~V.~Rashmi.
A Large Scale Analysis of Hundreds of In-Memory Cache Clusters at Twitter.
In \emph{USENIX Symposium on Operating Systems Design and Implementation (OSDI)}, 2020.

\bibitem{cachelib} B.~Berg, D.~S.~Berger, S.~McAllister, I.~Grosof, S.~Gunasekar, J.~Lu, M.~Uhlar, J.~Carrig, N.~Beckmann, M.~Harchol-Balter, and G.~R.~Ganger.
The CacheLib Caching Engine: Design and Experiences at Scale.
In \emph{USENIX Symposium on Operating Systems Design and Implementation (OSDI)}, 2020.

\bibitem{tinylfu} G.~Einziger, R.~Friedman, and B.~Manes.
TinyLFU: A Highly Efficient Cache Admission Policy.
\emph{ACM Transactions on Storage}, 13(4), 2017. arXiv:1512.00727.

\bibitem{slru} R.~Karedla, J.~S.~Love, and B.~G.~Wherry.
Caching Strategies to Improve Disk System Performance.
\emph{IEEE Computer}, 27(3):38--46, 1994.

\bibitem{clock} F.~J.~Corbat\'o.
A Paging Experiment with the Multics System.
In \emph{In Honor of Philip M. Morse} (H.~Feshbach and K.~U.~Ingard, eds.), MIT Press, 1969.

\bibitem{lsm} P.~O'Neil, E.~Cheng, D.~Gawlick, and E.~O'Neil.
The Log-Structured Merge-Tree (LSM-Tree).
\emph{Acta Informatica}, 33:351--385, 1996. DOI 10.1007/s002360050048.

\bibitem{rocksdb} Meta Platforms, Inc.
RocksDB: A Persistent Key-Value Store for Fast Storage Environments.
\url{https://rocksdb.org}, accessed 2026.

\bibitem{redis} Redis Ltd.
Redis: An In-Memory Data Store.
\url{https://redis.io}, accessed 2026.

\bibitem{iouring} J.~Axboe.
Efficient IO with io\_uring, 2019; and io\_uring(7), Linux Programmer's Manual, \url{https://man7.org/linux/man-pages/man7/io_uring.7.html}, accessed 2026.

\bibitem{splitmix} G.~L.~Steele~Jr., D.~Lea, and C.~H.~Flood.
Fast Splittable Pseudorandom Number Generators.
In \emph{ACM Conference on Object-Oriented Programming, Systems, Languages, and Applications (OOPSLA)}, pages 453--472, 2014. DOI 10.1145/2660193.2660195.

\bibitem{bitcask} J.~Sheehy and D.~Smith.
Bitcask: A Log-Structured Hash Table for Fast Key/Value Data.
Basho Technologies white paper, 2010. \url{https://riak.com/assets/bitcask-intro.pdf}.

\bibitem{silt} H.~Lim, B.~Fan, D.~G.~Andersen, and M.~Kaminsky.
SILT: A Memory-Efficient, High-Performance Key-Value Store.
In \emph{ACM Symposium on Operating Systems Principles (SOSP)}, pages 1--13, 2011. DOI 10.1145/2043556.2043558.

\bibitem{wisckey} L.~Lu, T.~S.~Pillai, A.~C.~Arpaci-Dusseau, and R.~H.~Arpaci-Dusseau.
WiscKey: Separating Keys from Values in SSD-Conscious Storage.
In \emph{USENIX Conference on File and Storage Technologies (FAST)}, 2016.

\bibitem{udepot} K.~Kourtis, N.~Ioannou, and I.~Koltsidas.
Reaping the Performance of Fast NVM Storage with uDepot.
In \emph{USENIX Conference on File and Storage Technologies (FAST)}, 2019.

\bibitem{s3fifo} J.~Yang, Y.~Zhang, Z.~Qiu, Y.~Yue, and R.~Vinayak.
FIFO Queues Are All You Need for Cache Eviction.
In \emph{ACM Symposium on Operating Systems Principles (SOSP)}, pages 130--149, 2023. DOI 10.1145/3600006.3613147.

\bibitem{sieve} Y.~Zhang, J.~Yang, Y.~Yue, Y.~Vigfusson, and K.~V.~Rashmi.
SIEVE is Simpler than LRU: An Efficient Turn-Key Eviction Algorithm for Web Caches.
In \emph{USENIX Symposium on Networked Systems Design and Implementation (NSDI)}, 2024.

\bibitem{fsyncfail} A.~Rebello, Y.~Patel, R.~Alagappan, A.~C.~Arpaci-Dusseau, and R.~H.~Arpaci-Dusseau.
Can Applications Recover from fsync Failures?
In \emph{USENIX Annual Technical Conference (ATC)}, 2020.

\bibitem{funsearch} B.~Romera-Paredes, M.~Barekatain, A.~Novikov, M.~Balog, M.~P.~Kumar, E.~Dupont, F.~J.~R.~Ruiz, J.~S.~Ellenberg, P.~Wang, O.~Fawzi, P.~Kohli, and A.~Fawzi.
Mathematical Discoveries from Program Search with Large Language Models.
\emph{Nature}, 625:468--475, 2024. DOI 10.1038/s41586-023-06924-6.

\bibitem{alphadev} D.~J.~Mankowitz, A.~Michi, A.~Zhernov, M.~Gelmi, M.~Selvi, C.~Paduraru, E.~Leurent, S.~Iqbal, J.-B.~Lespiau, A.~Ahern, T.~K\"oppe, K.~Millikin, S.~Gaffney, S.~Elster, J.~Broshear, C.~Gamble, K.~Milan, R.~Tung, M.~Hwang, A.~T.~Cemgil, M.~Barekatain, Y.~Li, A.~Mandhane, T.~Hubert, J.~Schrittwieser, D.~Hassabis, P.~Kohli, M.~Riedmiller, O.~Vinyals, and D.~Silver.
Faster Sorting Algorithms Discovered Using Deep Reinforcement Learning.
\emph{Nature}, 618(7964):257--263, 2023. DOI 10.1038/s41586-023-06004-9.

\bibitem{rewardhacking} J.~Skalse, N.~H.~R.~Howe, D.~Krasheninnikov, and D.~Krueger.
Defining and Characterizing Reward Hacking.
arXiv:2209.13085, 2022.

\bibitem{flashield} A.~Eisenman, A.~Cidon, E.~Pergament, O.~Haimovich, R.~Stutsman, M.~Alizadeh, and S.~Katti.
Flashield: A Hybrid Key-Value Cache that Controls Flash Write Amplification.
In \emph{USENIX Symposium on Networked Systems Design and Implementation (NSDI)}, pages 65--78, 2019.

\bibitem{ftwo} K.~Kanellis, B.~Chandramouli, T.~Hart, and S.~Venkataraman.
From FASTER to F2: Evolving Concurrent Key-Value Store Designs for Large Skewed Workloads.
\emph{Proceedings of the VLDB Endowment}, 18(12):4910--4923, 2025. DOI 10.14778/3750601.3750615.

\bibitem{kisssorcar} K.~Sen.
KISS Sorcar: A Stupidly-Simple General-Purpose and Software Engineering AI Assistant.
\url{https://github.com/ksenxx/kiss_ai} (paper at \texttt{papers/kisssorcar/}), 2026.

\bibitem{alphaevolve} A.~Novikov, N.~V\~u, M.~Eisenberger, E.~Dupont, P.-S.~Huang, A.~Z.~Wagner, S.~Shirobokov, B.~Kozlovskii, F.~J.~R.~Ruiz, A.~Mehrabian, M.~P.~Kumar, A.~See, S.~Chaudhuri, G.~Holland, A.~Davies, S.~Nowozin, P.~Kohli, and M.~Balog.
AlphaEvolve: A Coding Agent for Scientific and Algorithmic Discovery.
arXiv:2506.13131, 2025.

\bibitem{gepa} L.~A.~Agrawal, S.~Tan, D.~Soylu, N.~Ziems, R.~Khare, K.~Opsahl-Ong, A.~Singhvi, H.~Shandilya, M.~J.~Ryan, M.~Jiang, C.~Potts, K.~Sen, A.~G.~Dimakis, I.~Stoica, D.~Klein, M.~Zaharia, and O.~Khattab.
GEPA: Reflective Prompt Evolution Can Outperform Reinforcement Learning.
In \emph{International Conference on Learning Representations (ICLR)}, 2026. arXiv:2507.19457.

\end{thebibliography}

\end{document}
